\pagebreak

\chapter{Inbuilt Functions}

    \index{function!inbuilt}
    The following functions are built into the compiler. User functions must
    not use identifiers that conflict with these.
    
    The following abbreviations are used to denote modes -

    \begin{tabular}{| l | l |}
    \hline
    I  & immediate                      \\ \hline
    V  & value                          \\ \hline
    SV & selectvalue                    \\ \hline
    S  & static                         \\ \hline
    Q  & queue                          \\ \hline
    CM & combinatorial output memory    \\ \hline
    RM & registered output memory       \\ \hline
    PR & priority                       \\ \hline
    IN & input                          \\ \hline
    O  & output                         \\ \hline
    C  & clock                          \\ \hline
    \end{tabular}
    
    The following abbreviations are used to denote types -

    \begin{tabular}{| l | l |}
    \hline
    bits  & bits                                                    \\ \hline
    uint  & unsigned integer                                        \\ \hline
    int   & integer or unsigned integer                             \\ \hline
    fixed & fixed point                                             \\ \hline
    float & floating point                                          \\ \hline
    log   & logical (boolean)                                       \\ \hline
    str   & string                                                  \\ \hline
    enum  & enumerated type                                         \\ \hline
    type  & type                                                    \\ \hline
    file  & file                                                    \\ \hline
    map   & map                                                     \\ \hline
    list  & list                                                    \\ \hline
    any   & any type                                                \\ \hline
    prim  & a primitive type                                        \\ \hline
    comp  & a compound type (arrays, structs or any combination)    \\ \hline
    []    & array of type                                           \\ \hline
    \end{tabular}

    Functions with 'yes' to 'Parameter=argument allowed' can have arguments
    matched to parameters by position in the argument list or by
    parameter=argument (keymatch) or a mixture as long as positional arguments
    do not follow keymatch arguments. The keymatch arguments may be in
    any order. Otherwise these keymatch arguments are
    not allowed. This may be because - \\
    \blist
    \item   the function has a single input or output argument in which case
    no positional or keymatch parameter association is needed
    \item   the function has an indefinite number of arguments and hence there
    are no 'named' parameters with which to associate arguments
    \blend


    \section{abs - absolute value function}\label{sec:ifabs}
        
        
        {\bf Call:} \\
        abs(v)
        
        {\bf parameters:} \\
        \begin{tabular}{| l | l | l | l |}
        \hline
        param  & mode & type & \\
        \hline \hline
        v & any & int, fixed or float & \\ \hline
        \end{tabular}
        
        {\bf Parameter=argument allowed:} \\
        No
        
        {\bf return mode:} \\
        immediate or target

        {\bf return type:} \\
        int, fixed or float as per argument
        
        {\bf Description:} \\
        Returns the absolute value of the argument. The mode of the input
        may be IMMEDIATE, VALUE, SELECTVALUE, STATIC or QUEUE.
        
        An IMMEDIATE mode input may only be type "int" or "float"
        and the result will be the same type.
        
        Target mode inputs may only be types "int:", "fixed:." or "float:."
        and the result will be the same type.


    \section{accept - ignore clock domain differences}\label{sec:ifaccept}
        
        {\bf Call:} \\
        accept(v, c)
        
        {\bf parameters:} \\
        \begin{tabular}{| l | l | l | l |}
        \hline
        param  & mode & type & \\
        \hline \hline
        v & V & any & variable or expression \\ \hline
        c & C &  -  & optional clock variable or null \\ \hline
        \end{tabular}
        
        {\bf Parameter=argument allowed:} \\
        No
        
        {\bf return mode:} \\
        value

        {\bf return type:} \\
        type of argument
        
        {\bf Description:} \\
        This function returns a value mode result which is the argument variable or
        expression with its clock domain changed. If no second argument is provided
        the returned value is the same as the first argument but with the current
        clock domain. If the second argument is provided it should be a clock mode
        variable identifier and the returned value is the same as the first
        argument but with the designated clock domain. If the second argument is
        {\bf null} or the comma is present but the argument is omitted, the
        returned value is the same as the first argument but with no clock domain.
        The clock domain of the argument is neither checked nor resolved
        (assigned).
        
        If the first argument clock domain is the same as the current clock domain
        the argument is simply returned as the value.        
        
        The first argument cannot contain queue reads and cannot be a memory or
        clock mode variable.
        
        This function is used where it is desired to use the value of a
        variable or expression in a clock domain that differs from that of the
        argument without synchronising it in any way.
        
        This function should only be used in the following circumstances -
        \blist
        \item   the frequencies of the clocks of the two clock domains have
                a simple integer relationship, the clocks are locked in phase
                and the lower frequency clock edges are coincident with
                an edge of the higher frequency clock
        \item   the argument variable or expression only changes occasionally
                and the occurrence of corruption on a small proportion
                of data is acceptable
        \blend
        \index{constraint!timing}
        \index{timing constraint}
        In the former case a time constraint from the first clock domain to
        the second clock domain should be imposed, the value being the
        shortest time interval between a clock edge in the source domain and
        the next edge in the destination domain. In the latter case a time
        constraint would be meaningless.
        
        If Clocks are not suitably related and corrupted data is not acceptable
        then the inbuilt function {\bf sample()} \autorefph{sec:ifsample} should be used
        instead.
        
        {\bf WARNING:} \\
        {\bf This function must be used with great care! It introduces the
        possibility that data may be changing at the time of the local clock
        transition, resulting in data corruption. The consequences can be
        serious if this occurs in a test value in a {\em when}, target {\em
        while} or target {\em do while} structure in which case execution
        control can be unpredictably and permanently disrupted\footnote{This can
        occur because the control logic usually clocks the single bit test value
        into two flip-flops simultaneously. The routing delays to the flip-flops
        will differ and if the signal is glitching during setup the flip-flops
        may be set inconsistently. This could result in failure to execute
        either part of a when statement (which will kill the whole execution
        thread) or execution of both parts. Even worse disruption could
        occur in the case of loops.}. The use of {\em accept()} in such
        contexts has not been prohibited since there are situations where it is
        safe. This is usually in the context of external asynchronous signals
        and requires the programmer to have (and correctly apply) detailed
        knowledge of signal timing.}


    \section{acos - trigonometric arc cosine function}\label{sec:ifacos}
        
        
        {\bf Call:} \\
        acos(f)
        
        {\bf parameters:} \\
        \begin{tabular}{| l | l | l | l |}
        \hline
        param  & mode & type & \\
        \hline \hline
        f & I & float & \\ \hline
        \end{tabular}
        
        {\bf Parameter=argument allowed:} \\
        No
        
        {\bf return mode:} \\
        immediate

        {\bf return type:} \\
        float
        
        {\bf Description:} \\
        Returns the arc cosine of the immediate argument. The result is the
        angle in radian.


    \subsection{addsub - conditionally add or subtract}\label{sec:ifaddsub}

        {\bf Call:} \\
        addsub(in1, in2, add, gate, xci)
        
        {\bf parameters:} \\
        \begin{tabular}{| l | l | l | l |}
        \hline
        param  & mode & type & \\
        \hline \hline
        in1  & V,SV,S,Q & uint, int & \\ \hline
        in2  & V,SV,S,Q & uint, int & \\ \hline
        add  & V,SV,S,Q & log       & \\ \hline
        gate & V,SV,S,Q & log       & \\ \hline
        xci  & V,SV,S,Q & log       & \\ \hline
        \end{tabular}
        
        {\bf Parameter=argument allowed:} \\
        Yes
        
        {\bf return mode:} \\
        value

        {\bf return type:} \\
        uint, int, ufixed, fixed
        
        {\bf Description:} \\
        This function returns the sum or difference of its operands. If
        the input parameter {\bf add} is true, the value returned is {\bf
        in1} + {\bf in2}, otherwise it is {\bf in1} - {\bf in2}.
        
        If the optional input parameter {\bf gate} has an argument then when the
        value of that argument is false the second input argument is gated to zero
        and thus the result is the first input argument regardless of whether the
        operation is add or subtract. When {\bf gate} is false both {\bf add}
        and {\bf xci} have no effect.
        
        If the optional input parameter {\bf xci} has an argument then it is XORed
        with the carry in. This has the effect of incrementing the sum in the case
        of an addition or decrementing the difference in the case of a subtraction.
        Using the carry in this way is much more efficient than using a second
        adder to increment the sum or a second subtracter to decrement the
        difference.
        
        The following table shows the result type for allowed input types.
        Types with a trailing ":" are target types.

        \begin{tabular}{| l || l | l | l | l | l | l |}
        \hline
                & uint    & int    & uint:   & int:   & ufixed: & fixed: \\
        \hline \hline
        uint    & uint    & int    & uint:   & int:   & ufixed: & fixed: \\ \hline
        int     & int     & int    & int:    & int:   & fixed:  & fixed: \\ \hline
        uint:   & uint:   & int:   & uint:   & int:   & ufixed: & fixed: \\ \hline
        int:    & int:    & int:   & int:    & int:   & fixed:  & fixed: \\ \hline
        ufixed: & ufixed: & fixed: & ufixed: & fixed: & ufixed: & fixed: \\ \hline
        fixed:  & fixed:  & fixed: & fixed:  & fixed: & fixed:  & fixed: \\ \hline
        \end{tabular}
        
        If the operands are mixed signed/unsigned then the unsigned operand is
        increased in width by one by appending a zero bit at the most significant
        end. Integer operands are matched in size by padding the smaller operand.
        If the operands are mixed integer/fixed point the operands are aligned at
        the fixed point and padded at either end to match their sizes.

        The returned value has the same number of significant bits as the maximum
        of the widths of {\bf in1} and {\bf in2} plus one (extra MSB).
        
        If one or both operands are fixed point then the operand with the smallest
        fixed point offset is left shifted to align the binary points prior to
        adding or subtracting. The result width is the maximum of the aligned
        operand widths plus one (extra MSB) and the result binary point offset is
        the same as the aligned operand binary point offset.
        
        Any or all operands may contain queue reads, in which case the availability
        of the function value will depend on the availability of the operands
        which contain queue reads.


    \section{arglist - pass lists, arrays or maps of arguments}\label{sec:ifarglist}
        
        
        {\bf Call:} \\
        arglist(k, v) \\
        or \\
        arglist(m)
        
        {\bf parameters:} \\
        \begin{tabular}{| l | l | l | l |}
        \hline
        param  & mode & type & \\
        \hline \hline
        k & I & array or list & \\ \hline
        v & I & array or list & \\ \hline
        m & I & map           & \\ \hline
        \end{tabular}
        
        {\bf Parameter=argument allowed:} \\
        No
        
        {\bf return mode:} \\
        immediate

        {\bf return type:} \\
        -
        
        {\bf Description:} \\
        This function allows lists or maps of arguments to be inserted at any point within an
        input or output argument list. When called with two arguments the first is an array or
        a list of key strings. This argument may be explicitly 'null' (i.e. not omitted) if
        there are no key=value style arguments. The second argument is an array or list of
        arguments to be inserted. If both arguments are present the two arrays/lists must be
        the same size. The key array/list may have leading null entries as long as the
        complete resulting argument list to the enclosing module, procedure or function does
        not have key=value arguments preceding normal arguments (this is not illegal but is
        very confusing and should be avoided).
        
        If this function is called with a single argument then it must be type "map"
        The map entries are inserted as key=value arguments.
        
        Amongst other uses this function can be used to pass varargs arguments to
        deeper subroutines, e.g.
        
        \begin{lstlisting}[gobble=12, language=threepl]
            procedure p(varargs <- varargs) {
                p2(arglist(null, outargs) <- arglist(null, inargs));
            }
        \end{lstlisting}
        
        or
        
        \begin{lstlisting}[gobble=12, language=threepl]
            procedure p(varargs <- varargs) {
                p2(arglist(outkeys, outargs) <- arglist(inkeys, inargs));
            }
        \end{lstlisting}
        
        or more generally arguments from a procedure call may be inserted
        amongst other arguments, e.g.

        
        \begin{lstlisting}[gobble=12, language=threepl]
            procedure p(varargs) {
                .
                .
                p2(arglist(a, b, arglist(null, inargs), c);
                .
            }
        \end{lstlisting}
        
        
        This function should only be called as an argument to a library or
        user-defined class, module, procedure or function. It should not be
        called as an argument to an inbuilt class, module, procedure or
        function. In that case it will result in error termination.


    \section{arraytostr - convert an array of integers to a string}\label{sec:ifarraytostr}
        
        {\bf Call:} \\
        arraytostr(ia)
        
        {\bf parameters:} \\
        \begin{tabular}{| l | l | l | l |}
        \hline
        param  & mode & type & \\
        \hline \hline
        ia & I & []uint & integer array   \\ \hline
        \end{tabular}
        
        {\bf Parameter=argument allowed:} \\
        No
        
        {\bf return mode:} \\
        immediate

        {\bf return type:} \\
        \verb'str'
        
        {\bf Description:} \\
        Converts an array of integers, whose values are ASCII codes,
        into a string.

        See also - \\
        \autorefph{sec:ifstrtoarray} {\bf strtoarray()} converting
        a string to an integer array.


    \section{arraytype - array type of a variable or expression}\label{sec:ifarraytype}
        
        {\bf Call:} \\
        arraytype(v\_or\_e)
        
        {\bf parameters:} \\
        \begin{tabular}{| l | l | l | l |}
        \hline
        param  & mode & type & \\
        \hline \hline
        v\_or\_e & - &  -  & variable or expression \\ \hline
        \end{tabular}
        
        {\bf Parameter=argument allowed:} \\
        No
        
        {\bf return mode:} \\
        immediate

        {\bf return type:} \\
        type
        
        {\bf Description:} \\
        Return the array type of a variable or expression.
        If the type is not an array null is returned.
        
        See also - \\
        {\bf type()} \autorefph{sec:iftype} \\
        {\bf typestring()} \autorefph{sec:iftypestring} \\
        {\bf primtypestring()} \autorefph{sec:ifprimtypestring} \\
        {\bf imtotargtype()} \autorefph{sec:ifimtotargtype} \\
        {\bf targtoimtype()} \autorefph{sec:iftargtoimtype} \\
        {\bf typeval()} \autorefph{sec:iftypeval} \\
        {\bf typevalstring()} \autorefph{sec:iftypevalstr} \\
        {\bf typeisequal()} \autorefph{sec:iftypeiseq} \\


    \section{asin - trigonometric arc sine function}\label{sec:ifasin}
        
        
        {\bf Call:} \\
        asin(f)
        
        {\bf parameters:} \\
        \begin{tabular}{| l | l | l | l |}
        \hline
        param  & mode & type & \\
        \hline \hline
        f & I & float & \\ \hline
        \end{tabular}
        
        {\bf Parameter=argument allowed:} \\
        No
        
        {\bf return mode:} \\
        immediate

        {\bf return type:} \\
        float
        
        {\bf Description:} \\
        Returns the arc sine of the immediate argument. The result is the
        angle in radian.


    \section{atan - single argument trigonometric arc tangent function}\label{sec:ifatan}
        
        
        {\bf Call:} \\
        atan(f)
        
        {\bf parameters:} \\
        \begin{tabular}{| l | l | l | l |}
        \hline
        param  & mode & type & \\
        \hline \hline
        f & I & float & \\ \hline
        \end{tabular}
        
        {\bf Parameter=argument allowed:} \\
        No
        
        {\bf return mode:} \\
        immediate

        {\bf return type:} \\
        float
        
        {\bf Description:} \\
        Returns the arc tangent of the immediate argument. The result is the
        angle in radian.


    \section{atan2 - two argument trigonometric arc tangent function}\label{sec:ifatan2}
        
        
        {\bf Call:} \\
        atan2(f1, f2)
        
        {\bf parameters:} \\
        \begin{tabular}{| l | l | l | l |}
        \hline
        param  & mode & type & \\
        \hline \hline
        f1 & I & float & \\ \hline
        f2 & I & float & \\ \hline
        \end{tabular}
        
        {\bf Parameter=argument allowed:} \\
        No
        
        {\bf return mode:} \\
        immediate

        {\bf return type:} \\
        float
        
        {\bf Description:} \\
        Returns the arc tangent of the 1st argument divided by the 2nd
        argument. The arguments must both be immediate mode. The result is the
        angle in radian.


    \section{attribute - get an attribute from a variable}\label{sec:ifattribute}
        
        
        {\bf Call:} \\
        attribute(v, attr)
        
        {\bf parameters:} \\
        \begin{tabular}{| l | l | l | l |}
        \hline
        param  & mode & type & \\
        \hline \hline
        v        & I & any & \\ \hline
        attr     & I & str & \\ \hline
        \end{tabular}
        
        {\bf Parameter=argument allowed:} \\
        No
        
        {\bf return mode:} \\
        immediate

        {\bf return type:} \\
        type of the attribute
        
        {\bf Description:} \\
        Returns an attribute from the attributes map for the variable. For a
        list of attributes see \autorefph{sec:attributes}. The function will
        return attributes which are in the extended category.
        
        The first argument is the variable identifier.
        
        The 2nd argument is a string giving the attribute identifier.
        
        If the specified attribute is not present in the attribute map then
        null is returned.


    \section{attributes - get the attributes of a variable}\label{sec:ifattributes}
        
        
        {\bf Call:} \\
        attributes(v, extended)
        
        {\bf parameters:} \\
        \begin{tabular}{| l | l | l | l |}
        \hline
        param  & mode & type & \\
        \hline \hline
        v        & I & any & \\ \hline
        extended & I & log & \\ \hline
        \end{tabular}
        
        {\bf Parameter=argument allowed:} \\
        No
        
        {\bf return mode:} \\
        immediate

        {\bf return type:} \\
        float
        
        {\bf Description:} \\
        Returns the attributes map for the variable. The returned map is not a
        copy (except when 'extended' is true) so any changes o the map will be
        reflected in the map withion the variable itself.
        For a list of attributes see \autorefph{sec:attributes}.
        
        If the 2nd argument is present and true, some internal attributes
        of the variable will also be returned. These extra attributes are -
        \blist
        \item    {\bf identifier} - the abbreviated identifier
        \item    {\bf eidentifier} - the extended identifier
        \item    {\bf mode} - the mode
        \item    {\bf type} - the type
        \item    {\bf typestring} - the type as a string
        \item    {\bf decloc} - the source file location where created
        \item    {\bf inputparam} - true if this is an input parameter
        \item    {\bf outputparam} - true if this is an output parameter
        \item    {\bf matched} - true if this is a parameter matched by an argument
        \item    {\bf readonly} - true if read-only
        \item    {\bf assigned} - true if the variable has been assigned
        \item    {\bf used} - true if it is a target mode variable and its
                            value has been used
        \item    {\bf indirect} - true if the variable is indirect. This is the case
                            for some parameter/argument situations and is also the
                            case for clocks assigned to other clocks or clocks
                            which are inverted.
        \item    {\bf clk\_sinks} - if it is a clock mode variable the
                            number of destinations to which the clock is connected
        \item    {\bf clk\_srcs} - if it is a clock mode variable the
                            number of sources from which the clock is driven (should be 1)
        \blend
        
        Note that these attributes cannot be set so this returned map should not
        be applied to a variable. If the second argument is omitted or false
        the extended attributes will be omitted and the returned attributes map
        is suitable for application to a variable.
        
        Note also that an extended attributes map is a new map containing attributes
        copied from the attribute map of the variable. Any changes to the map will not
        affect the variable attributes.


    \section{biasedexponent - target float biased exponent}\label{sec:ifbiasedexponent}
        
        {\bf Call:} \\
        biasedexponent(arg)
        
        {\bf parameters:} \\
        \begin{tabular}{| l | l | l | l |}
        \hline
        param  & mode & type & \\
        \hline \hline
        arg  & V SV S Q IN & float & target float value  \\ \hline
        \end{tabular}
        
        {\bf Parameter=argument allowed:} \\
        No
        
        {\bf return mode:} \\
        value

        {\bf return type:} \\
        uint
        
        {\bf Description:} \\
        This function returns the biased exponent of a target floating point
        value. The biased exponent is return as type uint: whose width is
        equal to the width of the biased exponent.


    \section{biasedexponentwidth - width of target float biased exponent}\label{sec:ifbiasedexponentwidth}
        
        {\bf Call:} \\
        biasedexponentwidth(arg)
        
        {\bf parameters:} \\
        \begin{tabular}{| l | l | l | l |}
        \hline
        param  & mode & type & \\
        \hline \hline
        arg  & I V SV S Q IN & type float & target float type or value  \\ \hline
        \end{tabular}
        
        {\bf Parameter=argument allowed:} \\
        No
        
        {\bf return mode:} \\
        immediate

        {\bf return type:} \\
        uint
        
        {\bf Description:} \\
        This function returns the number of bits in the biased exponent of
        a target floating point value. The argument may be a target floating
        point value or it may be a target floating point type.


    \section{bitsfor - number of bits to represent an integer - DEPRECATED!}\label{sec:ifbitsfor}
        
        {\bf Call:} \\
        bitsfor(arg)
        
        {\bf parameters:} \\
        \begin{tabular}{| l | l | l | l |}
        \hline
        param  & mode & type & \\
        \hline \hline
        arg  & I & uint or int & integer value  \\ \hline
        \end{tabular}
        
        {\bf Parameter=argument allowed:} \\
        No
        
        {\bf return mode:} \\
        immediate

        {\bf return type:} \\
        uint
        
        {\bf Description:} \\
        This function returns the number of bits needed to represent the +ve or -ve
        integer argument. For a +ve integer the sign bit is not included.
        
        The function {\em size()} \autorefph{sec:ifsize} should be used instead. It
        will return the same value, and is more general in that it will accept
        target primitives as well as immediate arguments.


    \section{cast - cast a value to a another type or type width}\label{sec:ifcast}
        
        {\bf Call:} \\
        cast(value, type)
        
        {\bf parameters:} \\
        \begin{tabular}{| l | l | l | l |}
        \hline
        param  & mode & type & \\
        \hline \hline
        value & I, V, SV, S, Q & any       & value to be cast    \\ \hline
        type  & I              & str, type & type string or type \\ \hline
        \end{tabular}
        
        {\bf Parameter=argument allowed:} \\
        No
        
        {\bf return mode:} \\
        mode of first argument

        {\bf return type:} \\
        type of second argument
        
        {\bf Description:} \\
        For immediate types conversions between types {\em bits}, {\em uint},
        {\em int} and {\em str} are allowed with the provisos -
        \blist
        \item   int to uint gives an error if the integer is -ve
        \item   str to numeric gives an error if not correctly formatted
        \blend
        A string to be converted to numeric must be
        \blist
        \item   a string of 0's and 1's preceded by "0b" or 0B"
        \item   a string of hexadecimal digits preceded by "0x" or "0X"
        \item   a string of octal digits preceded by "0"
        \item   a string of decimal digits whose 1st digit is not 0
        \blend


        Target mode type conversions allowed are shown in the following
        table -

        \begin{tabular}{| l | l || l | l | l | l | l | l | l | l |}
        \hline
        \multicolumn{2}{| c |}{} & \multicolumn{8}{| c |}{to} \\
        \multicolumn{2}{| c |}{} & bits & uint & int & ufixed & fixed & log & float & enum \\
        \hline \hline
                        & bits   & az   & az   & as  & az     & as    & az  & az    & -  \\ \cline{2-10}
                        & uint   & az   & az   & az  & cz     & cz    & az  & -     & az \\ \cline{2-10}
                        & int    & as   & as   & as  & cs     & cs    & -   & -     & -  \\ \cline{2-10}
                   from & ufixed & az   & bz   & bz  & dz     & dz    & -   & -     & -  \\ \cline{2-10}
                        & fixed  & as   & -    & bs  & ds     & ds    & -   & -     & -  \\ \cline{2-10}
                        & log    & az   & az   & -   & -      & -     & a   & -     & -  \\ \cline{2-10}
                        & float  & az   & -    & -   & -      & -     & -   & e     & -  \\ \cline{2-10}
                        & enum   & -    & az   & -   & -      & -     & -   & -     & -  \\
        \hline
        \end{tabular}

        \blist
        \item[{\bf a}] copy right justified
        \item[{\bf b}] copy from left of binary point, ignoring fractional part
        \item[{\bf c}] copy to left of binary point and zero fractional part
        \item[{\bf d}] copy aligned with binary point truncating or zero padding least significant
          bits if size of fraction differs
        \item[{\bf e}] copy with adjustments to biased exponent and mantissa as required
        \blend

        \blist
        \item[{\bf z}] zero pad or truncate most significant end if width differs
        \item[{\bf s}] sign extend or truncate most significant end if width differs
        \blend

        If casting to or from compound types any types are allowed (including
        primitive types) as long as the numbers of bits are identical.
        Conversions into or out of compound types is carried out by mapping
        the bits in order without regard to component boundaries.

        When casting between int or uint and fixed or ufixed the binary
        point is recognised and where appropriate the sign is
        propagated, i.e. the cast retains numeric sense.

        When casting between bits and fixed or ufixed the binary point
        is not recognised and bits are simply right justified.

        When casting between bits and int or fixed and the cast type is
        wider, sign extension is performed instead of zero padding.

        Type float can only be cast to another float, adjusting the
        biased exponent and mantissa as required.

        Type bits or uint can be cast to type log which will extract the
        rightmost bit as false (0) or true (1). Casting type log to bits
        or uint results in a single bit, zero padded where required.

        Casting an enum to a uint provides the ordinal value (there is
        also an ord() built in function). Casting a uint to an enum
        selects the enum value given the ordinal.
        
        As a special case a queue of type "null" can be cast to a value of
        type "log".
        
        See also - \\
        {\bf round()} \autorefph{sec:ifround}) \\
        {\bf trunc()} \autorefph{sec:iftrunc}) \\
        {\bf ceil()} \autorefph{sec:ifceil}) \\
        {\bf floor()} \autorefph{sec:iffloor}) \\
        {\bf ord()} \autorefph{sec:iford})
        


    \section{ceil - round float to nearest integer towards $+\infty$}\label{sec:ifceil}
        
        
        {\bf Call:} \\
        ceil(f)
        
        {\bf parameters:} \\
        \begin{tabular}{| l | l | l | l |}
        \hline
        param  & mode & type & \\
        \hline \hline
        f & I & float & \\ \hline
        \end{tabular}
        
        {\bf Parameter=argument allowed:} \\
        No
        
        {\bf return mode:} \\
        immediate

        {\bf return type:} \\
        int
        
        {\bf Description:} \\
        Rounds the floating point argument to the nearest integer towards
        $+\infty$.
        An integer type is returned


    \section{cmemoryread - read a combinatorial output memory}\label{sec:ifcmemoryread}
        
        {\bf Call:} \\
        cmemoryread(id, port, address)
        
        {\bf parameters:} \\
        \begin{tabular}{| l | l | l | l |}
        \hline
        param  & mode & type & \\
        \hline \hline
        id      & M        &  -  & memory identifier    \\ \hline
        port    & I        & int & port                 \\ \hline
        address & V SV S Q & any & address value        \\ \hline
        \end{tabular}
        
        {\bf Parameter=argument allowed:} \\
        Yes
        
        {\bf return mode:} \\
        value

        {\bf return type:} \\
        data type of memory
        
        {\bf Description:} \\
        {\em id} is the memory mode variable identifier.
        
        {\em port} is the memory port.
        
        {\em address} is the address.
        
        The address may be any type, including a compound type, but the type
        must match the type given when the registered output memory
        was created (see {\bf cmemory()} \autorefph{sec:ipcmemory}) but primitive component
        widths may differ in which case padding or truncation will occur.
             
        The value of the function is the data output from the read. The type of
        the function return is the same type as the data type given when the
        combinatorial output memory was created.
        
        There may be calls to read and/or write a memory port at multiple
        points in the program, but the program logic must ensure that these
        calls are mutually exclusive (the compiler cannot determine that).
        Failure to observe this will result in data errors.
        
        If both the following conditions hold -
        \blist
        \item   a port has only a single call to cmemoryread and no calls to
                procedure cmemorywrite()
        \item   the address expression contains no queues
        \blend
        then the function output will continuously (statically) reflect the
        memory contents at the requested address. This is because for a single
        read access only the address input is permanently connected and hence
        the output remains continuously valid. If the port has more than one
        access point then the address must be multiplexed and is only
        connected for the execution cycle of the statement in which the
        function call appears. For a single read access expressions containing
        a call to this function can be passed as arguments to modules,
        procedures or functions or can be assigned to value variables.

        If the above is not the case then the function output will only be
        valid in the execution cycle of the assignment statement in which the
        function call is part of the RHS expression. In this case an expression
        containing {\bf cmemoryread()} cannot be assigned to a value variable
        after that value has been evaluated (this is detected and results in
        a fatal error).
        
        A call to {\bf cmemorywrite()} (\autorefph{sec:ipcmemorywrite}) or {\bf cmemoryread()}
        (\autorefph{sec:ifcmemoryread}) sets the clock domain for the memory port to the
        current clock domain. All further calls to either of these inbuilt
        procedures/functions for the same port of the same combinatorial output
        memory must be in the same clock domain (or a fatal error results).
            
        Note that {\bf cmemoryread()} may appear in an expression assigned
        to a value mode variable and that variable may then be used in the
        RHS of a subsequent assignment statement. However, there are two limitations
        to the use of {\bf cmemoryread()} in this context -
        \blist
        \item   the value variable must be used at least once in the RHS of
                an assignment otherwise the control logic for the memory
                read will have no inputs and a code generation error
                will occur
        \item   the value variable must be assigned the expression
                containing {\bf cmemoryread()} prior to its occurrence
                in the RHS of an assignment - post-assignment cannot
                be used in this instance
        \blend
        These errors are not recognised during immediate execution but
        produce indecipherable errors during code generation
        (proper execution checks will be implemented later).




%    \section{componentpointer - get a pointer to a variable component}\label{sec:ifcomponentpointer}
%        
%        {\bf Call:} \\
%        componentpointer(v, i)
%        
%        {\bf parameters:} \\
%        \begin{tabular}{| l | l | l | l |}
%        \hline
%        param  & mode & type & \\
%        \hline \hline
%        v & V S Q & any & static variable   \\ \hline
%        i & I & int or uint & index     \\ \hline
%        \end{tabular}
%        
%        {\bf Parameter=argument allowed:} \\
%        No
%        
%        {\bf return mode:} \\
%        immediate
%
%        {\bf return type:} \\
%        \verb'->'
%        
%        {\bf Description:} \\
%        Returns a pointer to a component of a value, static or queue mode
%        variable. This is provided to allow access to a single component of a
%        variable of compound type which has nested structs and/or arrays by
%        means of a single integer index instead of field names and arrays
%        subscripts. It has two arguments. The first argument is a variable
%        reference, i.e. it may have fields or subscripts. The second argument is
%        an integer index. If the variable type is primitive the index must be
%        zero. The index must be non-negative and less than the number of
%        components in the type.
%        
%        The increment and decrement operators cannot be used on a dereferenced
%        pointer returned from componentpointer().
%        
%        This function is a bit obscure but is intended for assigning or evaluating value, static or
%        queue mode variable fields or array members from outside the FPGA, e.g.
%        from an attached CPU. Requests to assign to variable fields or array
%        members can be transmitted to the FPGA using simple integer indices to
%        select via an array of pointers the components to be referenced or
%        assigned.
%        
%        When determining the index value to be used the following rules are
%        pertinent -
%        \blist
%        \item   struct fields are packed in the order they are listed in the
%                type
%        \item   arrays are packed with last subscript varying most rapidly
%        \blend
%        
%        See also {\bf componenttype()} \autorefph{sec:ifcomponenttype}.




%    \section{componenttype - get the type of a variable component}\label{sec:ifcomponenttype}
%        
%        {\bf Call:} \\
%        componenttype(v, i)
%        
%        {\bf parameters:} \\
%        \begin{tabular}{| l | l | l | l |}
%        \hline
%        param  & mode & type & \\
%        \hline \hline
%        v & S & any & static variable   \\ \hline
%        i & I & int or uint & index     \\ \hline
%        \end{tabular}
%        
%        {\bf Parameter=argument allowed:} \\
%        No
%        
%        {\bf return mode:} \\
%        immediate
%
%        {\bf return type:} \\
%        type
%        
%        {\bf Description:} \\
%        Returns the type of a component of a value, static or queue mode
%        variable. This is provided to allow access to a single component of a
%        variable of compound type which has nested structs and/or arrays by
%        means of a single integer index instead of field names and arrays
%        subscripts. It has two arguments. The first argument is a variable
%        reference, i.e. it may have fields or subscripts. The second argument is
%        an integer index. If the variable type is primitive the index must be
%        zero. The index must be non-negative and less than the number of
%        components in the type.
%        
%        This function is a companion to inbuilt function
%        {\bf componentpointer()} \autorefph{sec:ifcomponentpointer}.



    \section{copy - copy a list or map}\label{sec:ifcopy}
        
        {\bf Call:} \\
        copy(v)
        
        {\bf parameters:} \\
        \begin{tabular}{| l | l | l | l |}
        \hline
        param  & mode & type & \\
        \hline \hline
        v & I & list or map & subject list or map   \\ \hline
        \end{tabular}
        
        {\bf Parameter=argument allowed:} \\
        No
        
        {\bf return mode:} \\
        immediate

        {\bf return type:} \\
        list or map
        
        {\bf Description:} \\
        Returns a {\bf shallow} copy of the argument list or map. Most immediate
        assignments are by value, i.e. the RHS values are copied to the LHS, but
        lists and maps are not copied on assignment. Following assignment the
        LHS and RHS reference the same object. This function is provided for
        occasions where a copy is required.


    \section{cos - trigonometric cosine function}\label{sec:ifcosin}
        
        
        {\bf Call:} \\
        cos(f)
        
        {\bf parameters:} \\
        \begin{tabular}{| l | l | l | l |}
        \hline
        param  & mode & type & \\
        \hline \hline
        f & I & float & \\ \hline
        \end{tabular}
        
        {\bf Parameter=argument allowed:} \\
        No
        
        {\bf return mode:} \\
        immediate

        {\bf return type:} \\
        float
        
        {\bf Description:} \\
        Returns the cosine of the immediate argument. The argument is the
        angle in radian.


    \section{cspan - span a head string selected by characters not in a set}\label{sec:ifcspan}
        
        {\bf Call:} \\
        cspan(s1, s2)
        
        {\bf parameters:} \\
        \begin{tabular}{| l | l | l | l |}
        \hline
        param  & mode & type & \\
        \hline \hline
        s1 & I & str & subject string   \\ \hline
        s2 & I & str & search string    \\ \hline
        \end{tabular}
        
        {\bf Parameter=argument allowed:} \\
        No
        
        {\bf return mode:} \\
        immediate

        {\bf return type:} \\
        int
        
        {\bf Description:} \\
        Returns the length of the head string from the first argument whose characters are
        {\bf not} contained in the set of characters given by the second
        argument string.


    \section{currentclock - get the current clock variable}\label{sec:ifcurrentclock}
        
        {\bf Call:} \\
        currentclock()
        
        {\bf parameters:} \\
        none
        
        {\bf return mode:} \\
        clock

        {\bf return type:} \\
        none
        
        {\bf Description:} \\
        Gets the current clock variable. If there is no current clock
        a fatal error occurs. Null will be returned if there
        is no current clock domain or the clock domain is explicitly
        null.



    \section{currentelement - get the current element identifier}\label{sec:ifcurrentelement}
        
        {\bf Call:} \\
        currentelement()
        
        {\bf parameters:} \\
        none
        
        {\bf return mode:} \\
        immediate

        {\bf return type:} \\
        str
        
        {\bf Description:} \\
        Gets the identifier of the current element. The current element is that
        element currently being constructed, or that was most recently
        constructed, using the inbuilt procedure  {\bf element()}
        \autorefph{sec:ipelement}. Null will be returned if that procedure has
        not yet been called.


    \section{deflate - compress a string}\label{sec:ifdeflate}
        
        {\bf Call:} \\
        deflate(s)
        
        {\bf parameters:} \\
        \begin{tabular}{| l | l | l | l |}
        \hline
        param  & mode & type & \\
        \hline \hline
        s & - & -   & string to be compressed \\ \hline
        \end{tabular}
        
        {\bf return mode:} \\
        immediate

        {\bf return type:} \\
        \verb'[]uint'
        
        {\bf Description:} \\
        Returns the compressed string as an array of unsigned 32-bit integers.
        This function was specifically created for storing compressed .info text
        in rmemory (block RAM) so that it can later be read by the CPU and
        uncompressed (inflated). It avoids the necessity of running gzip in a
        shell from within 3PL and then reading the resulting file, thus greatly
        simplifying the code.

        The compressed bytes are packed as 32-bit unsigned integers such that
        reading them into CPU memory will result in the correct byte order.

        The first word is a checksum of the remainder of the block. The second
        word is a count of the number of bytes of uncompressed data and the
        third word a count of the number of compressed bytes. The compressed
        data that follows these counts contains a 2-byte header relevant to zlib
        which is 0x78 (compression method) and 0x9C (flags). The last word is
        zero-padded as required.

        It would be preferable for the compression output to be in an array of
        byte, but 3PL does not have this type. It would not be appropriate for
        the output to be a string type.


    \section{dimension - get a dimension of an array or size of a map or list}\label{sec:ifdimension}
        
        {\bf Call:} \\
        dimension(ident, index)
        
        {\bf parameters:} \\
        \begin{tabular}{| l | l | l | l |}
        \hline
        param  & mode & type & \\
        \hline \hline
        ident & - & -   & variable identifier       \\ \hline
        index & I & int & optional dimension index  \\ \hline
        \end{tabular}
        
        {\bf Parameter=argument allowed:} \\
        No
        
        {\bf return mode:} \\
        immediate

        {\bf return type:} \\
        int
        
        {\bf Description:} \\
        Returns one dimension of an array variable. The optional second
        argument selects which dimension to report. If the 2nd argument is
        omitted the first or only dimension will be reported. If the variable
        type is not an array, 0 is returned. If the index equals or exceeds
        the number of dimensions a fatal error occurs.
        
        It will also return the number of entries in a map or list. In that case
        the second argument must be omitted.


    \section{dimensions - get the number of dimensions of an array}\label{sec:ifdimensions}
        
        {\bf Call:} \\
        dimensions(ident)
        
        {\bf parameters:} \\
        \begin{tabular}{| l | l | l | l |}
        \hline
        param  & mode & type & \\
        \hline \hline
        ident & - & - & variable identifier \\ \hline
        \end{tabular}
        
        {\bf Parameter=argument allowed:} \\
        No
        
        {\bf return mode:} \\
        immediate

        {\bf return type:} \\
        int
        
        {\bf Description:} \\
        Returns the number of dimensions of a variable. If the variable type
        is not an array, 0 is returned.


    \section{directive - get the value of a directive}\label{sec:ifdirective}
        
        {\bf Call:} \\
        directive(key)
        
        {\bf parameters:} \\
        \begin{tabular}{| l | l | l | l |}
        \hline
        param  & mode & type & \\
        \hline \hline
        key & I & str & directive key \\ \hline
        \end{tabular}
        
        {\bf Parameter=argument allowed:} \\
        No
        
        {\bf return mode:} \\
        immediate

        {\bf return type:} \\
        primitive
        
        {\bf Description:} \\
        Returns the value of a directive. If the directive is not found a
        fatal error occurs (use {\em exists()} to check that a directive
        exists if unsure).

        The key cannot have a scope prepended - directives are all global.

        see also inbuilt functions - \\
        {\bf directives()} \autorefph{sec:ifdirectives}. \\
        {\bf direxists()} \autorefph{sec:ifdirexists}.\\
        and inbuilt procedures \\
        {\bf directive} \autorefph{sec:ipdirective}.\\
        {\bf removedirective} \autorefph{sec:ipremovedirective}.


    \section{directives - get the map of directives}\label{sec:ifdirectives}
        
        {\bf Call:} \\
        directives()
        
        {\bf parameters:} \\
        none
        
        {\bf return mode:} \\
        immediate

        {\bf return type:} \\
        map
        
        {\bf Description:} \\
        Returns a map of directives.

        see also inbuilt functions - \\
        {\bf directive()} \autorefph{sec:ifdirective}.\\
        {\bf direxists()} \autorefph{sec:ifdirexists}.\\
        and inbuilt procedures \\
        {\bf directive} \autorefph{sec:ipdirective}.\\
        {\bf removedirective} \autorefph{sec:ipremovedirective}.

    \section{direxists - determine if a directive exists}\label{sec:ifdirexists}
        
        {\bf Call:} \\
        direxists(key)
        
        {\bf parameters:} \\
        \begin{tabular}{| l | l | l | l |}
        \hline
        param  & mode & type & \\
        \hline \hline
        key & I &  str  & key string \\ \hline
        \end{tabular}
        
        {\bf Parameter=argument allowed:} \\
        No
        
        {\bf return mode:} \\
        immediate

        {\bf return type:} \\
        log
        
        {\bf Description:} \\
        Search currently accessible scopes and determine if a directive with the
        key specified by the argument string exists.
        
        The key cannot have a scope prepended - directives are all global.
        
        see also inbuilt functions - \\
        {\bf varexists()} \autorefph{sec:ifvarexists}. \\
        {\bf directives()} \autorefph{sec:ifdirectives}. \\
        and inbuilt procedures \\
        {\bf directive} \autorefph{sec:ipdirective}.\\
        {\bf removedirective} \autorefph{sec:ipremovedirective}.


    \section{domain - get the clock domain for a variable}\label{sec:ifdomain}
        
        {\bf Call:} \\
        domain(ident)
        
        {\bf parameters:} \\
        \begin{tabular}{| l | l | l | l |}
        \hline
        param  & mode & type & \\
        \hline \hline
        ident & - &  -  & variable identifier \\ \hline
        \end{tabular}
        
        {\bf Parameter=argument allowed:} \\
        No
        
        {\bf return mode:} \\
        clock

        {\bf return type:} \\
        none
        
        {\bf Description:} \\
        Returns the clock domain for a variable. If the variable has not had
        its value referenced or been assigned so far, null will be returned as
        the variable will not yet have had its clock determined. The variable
        may be selectvalue, value, static or priority mode. It cannot be queue,
        memory or clock mode.
        
        See also {\bf writedomain()} \autorefph{sec:ifwritedomain} which gets the write
        (input or assignment) clock and {\bf readdomain()} \autorefph{sec:ifreaddomain} which
        gets the read (output or evaluation) clock.
        
        For a queue the read and write clocks may be the same or may differ.
        For selectvalue, value, static and priority mode variables there is only
        one clock and both inbuilt functions will return that clock.


    \section{dutycycle - get the duty cycle of a clock variable}\label{sec:ifdutycycle}
        
        {\bf Call:} \\
        dutycycle(v)
        
        {\bf parameters:} \\
        \begin{tabular}{| l | l | l | l |}
        \hline
        param  & mode & type & \\
        \hline \hline
        v & C & - & clock variable \\ \hline
        \end{tabular}
        
        {\bf Parameter=argument allowed:} \\
        No
        
        {\bf return mode:} \\
        immediate

        {\bf return type:} \\
        float
        
        {\bf Description:} \\
        Returns the duty cycle of a clock variable. This is the proportion
        of time the clock is high.
        
        See also inbuilt functions - \\
        {\bf frequency()} \autorefph{sec:iffrequency} \\
        {\bf period()} \autorefph{sec:ifperiod}
 

    \section{elapsedtime - get the elapsed time}\label{sec:ifelapsedtime}
        
        {\bf Call:} \\
        elapsedtime(start)
        
        {\bf parameters:} \\
        \begin{tabular}{| l | l | l | l |}
        \hline
        param  & mode & type & \\
        \hline \hline
        start & I & log & 3PL start time   \\ \hline
        \end{tabular}
        
        {\bf Parameter=argument allowed:} \\
        No
        
        {\bf return mode:} \\
        immediate

        {\bf return type:} \\
        "float"
        
        {\bf Description:} \\
        Returns the elapsed time in second since the start of immediate
        execution(argument is missing or false) or the elapsed time since the start
        of 3PL compilation of the source code (argument is true).
        
        The time granularity and accuracy depends on the operating system.

        See also - \\
        \autorefph{sec:iftime} {\bf time()} for date and time elements and \\
        \autorefph{sec:iftimestring} {\bf timestring()} for date and time as a
        string.


    \section{endswith - a string is the tail of another}\label{sec:ifendswith}
        
        {\bf Call:} \\
        endswith(s1, s2)
        
        {\bf parameters:} \\
        \begin{tabular}{| l | l | l | l |}
        \hline
        param  & mode & type & \\
        \hline \hline
        s1 & I & str & subject string   \\ \hline
        s2 & I & str & test string      \\ \hline
        \end{tabular}
        
        {\bf Parameter=argument allowed:} \\
        No
        
        {\bf return mode:} \\
        immediate

        {\bf return type:} \\
        log
        
        {\bf Description:} \\
        Returns true if the 2nd argument string is the tail
        of the 1st argument string.


    \section{entries - get the number of entries in a list or map - DEPRECATED}\label{sec:ifentries}
        
        {\bf Call:} \\
        entries(ident)
        
        {\bf parameters:} \\
        \begin{tabular}{| l | l | l | l |}
        \hline
        param  & mode & type & \\
        \hline \hline
        ident & - & -   & list or map identifier \\ \hline
        \end{tabular}
        
        {\bf Parameter=argument allowed:} \\
        No
        
        {\bf return mode:} \\
        immediate

        {\bf return type:} \\
        int
        
        {\bf Description:} \\
        Returns the number of entries in a list or map.
        
        This function has been deprecated. {\bf dimension()} now
        includes the same function.


    \section{env - get the environment variables}\label{sec:ifenv}
        
        {\bf Call:} \\
        env()
        
        {\bf parameters:} \\
        None.
        
        {\bf return mode:} \\
        immediate

        {\bf return type:} \\
        map
        
        {\bf Description:} \\
        Returns the environment variables as a map. The map keys are the
        environment variable identifiers and the associated map values
        are the values of the variables. The values are strings.


    \section{exists - determine if an identifier is recognised - DEPRECATED!}\label{sec:ifexists}
        
        {\bf Call:} \\
        exists(ident)
        
        {\bf parameters:} \\
        \begin{tabular}{| l | l | l | l |}
        \hline
        param  & mode & type & \\
        \hline \hline
        ident & - &  -  & identifier \\ \hline
        \end{tabular}
        
        {\bf Parameter=argument allowed:} \\
        No
        
        {\bf return mode:} \\
        immediate

        {\bf return type:} \\
        int
        
        {\bf Description:} \\
        Search currently accessible scopes and determine if an identifier is
        in any maps, returning an integer code. Integer values are -

        \begin{tabular}{| l | l |}
        \hline
        0 & unknown     \\ \hline
        1 & variable    \\ \hline
        2 & directive   \\ \hline
        \end{tabular}

        In the case of directive maps the identifier is tried as a string key. If
        the identifier has {\bf global.}, {\bf file.}, {\bf calling.} or {\bf
        local.} prepended then the search will be restricted to the specified
        scope.
        
        This function has been deprecated as the use of an integer returned
        value is clumsy. The inbuilt functions {\bf varexists()}
        \autorefph{sec:ifvarexists} or {\bf direxists()}
        \autorefph{sec:ifdirexists} should be used instead to test explicitly
        for variables or directives.



    \section{exp - e to a power}\label{sec:ifexp}
        
        {\bf Call:} \\
        exp(f)
        
        {\bf parameters:} \\
        \begin{tabular}{| l | l | l | l |}
        \hline
        param  & mode & type & \\
        \hline \hline
        f & I & float & \\ \hline
        \end{tabular}
        
        {\bf Parameter=argument allowed:} \\
        No
        
        {\bf return mode:} \\
        immediate

        {\bf return type:} \\
        float
        
        {\bf Description:} \\
        Returns Euler's number {\em e} raised to the power of the immediate
        argument.



    \section{family - string identifying platform}\label{sec:iffamily}
        
        {\bf Call:} \\
        family()
        
        {\bf parameters:} \\
        
        None.
        
        {\bf return mode:} \\
        immediate

        {\bf return type:} \\
        str
        
        {\bf Description:} \\
        Returns a string identifying the platform, e.g. "xc5v" for
        Virtex5.
 


    \section{filecanread - check if a file is readable}\label{sec:iffilecanread}
        
        {\bf Call:} \\
        filecanread(f)
        
        {\bf parameters:} \\
        \begin{tabular}{| l | l | l | l |}
        \hline
        param  & mode & type & \\
        \hline \hline
        f & I & str or file & file \\ \hline
        \end{tabular}
        
        {\bf Parameter=argument allowed:} \\
        No
        
        {\bf return mode:} \\
        immediate

        {\bf return type:} \\
        log
        
        {\bf Description:} \\
        Return true if file exists and is readable.
        The argument must be immediate mode and may be type "str"
        of "file". The former is a path name string, relative or absolute.
        The latter is a file variable specifying the file.
 


    \section{filecanwrite - check if a file is writable}\label{sec:iffilecanwrite}
        
        {\bf Call:} \\
        filecanwrite(f)
        
        {\bf parameters:} \\
        \begin{tabular}{| l | l | l | l |}
        \hline
        param  & mode & type & \\
        \hline \hline
        f & I & str or file & file \\ \hline
        \end{tabular}
        
        {\bf Parameter=argument allowed:} \\
        No
        
        {\bf return mode:} \\
        immediate

        {\bf return type:} \\
        log
        
        {\bf Description:} \\
        Return true if file exists and is writable.
        The argument must be immediate mode and may be type "str"
        of "file". The former is a path name string, relative or absolute.
        The latter is a file variable specifying the file.
 


    \section{filedelete - delete a file}\label{sec:iffiledelete}
        
        {\bf Call:} \\
        filedelete(f)
        
        {\bf parameters:} \\
        \begin{tabular}{| l | l | l | l |}
        \hline
        param  & mode & type & \\
        \hline \hline
        f & I & str & file name \\ \hline
        \end{tabular}
        
        {\bf Parameter=argument allowed:} \\
        No
        
        {\bf return mode:} \\
        immediate

        {\bf return type:} \\
        log
        
        {\bf Description:} \\
        Delete a file. Return true if successful.
 


    \section{fileexists - check if a file exists}\label{sec:iffileexists}
        
        {\bf Call:} \\
        fileexists(f)
        
        {\bf parameters:} \\
        \begin{tabular}{| l | l | l | l |}
        \hline
        param  & mode & type & \\
        \hline \hline
        f & I & str or file & file \\ \hline
        \end{tabular}
        
        {\bf Parameter=argument allowed:} \\
        No
        
        {\bf return mode:} \\
        immediate

        {\bf return type:} \\
        log
        
        {\bf Description:} \\
        Return true if the file exists.
        The argument must be immediate mode and may be type "str"
        of "file". The former is a path name string, relative or absolute.
        The latter is a file variable specifying the file.
 


    \section{filegetparent - get the path to the file}\label{sec:iffilegetparent}
        
        {\bf Call:} \\
        filegetparent(f)
        
        {\bf parameters:} \\
        \begin{tabular}{| l | l | l | l |}
        \hline
        param  & mode & type & \\
        \hline \hline
        f & I & str & file name \\ \hline
        \end{tabular}
        
        {\bf Parameter=argument allowed:} \\
        No
        
        {\bf return mode:} \\
        immediate

        {\bf return type:} \\
        str
        
        {\bf Description:} \\
        Get the path to the file. Return the path string.
 


    \section{fileisdir - check if a file is a directory}\label{sec:iffileisdir}
        
        {\bf Call:} \\
        fileisdir(f)
        
        {\bf parameters:} \\
        \begin{tabular}{| l | l | l | l |}
        \hline
        param  & mode & type & \\
        \hline \hline
        f & I & str or file & file \\ \hline
        \end{tabular}
        
        {\bf Parameter=argument allowed:} \\
        No
        
        {\bf return mode:} \\
        immediate

        {\bf return type:} \\
        log
        
        {\bf Description:} \\
        Return true if the file is a directory. The argument must be immediate
        mode and may be type "str" of "file". The former is a path name string,
        relative or absolute. The latter is a file variable specifying the
        file.
 


    \section{fileisfile - check if a file is a file}\label{sec:iffileisfile}
        
        {\bf Call:} \\
        fileisfile(f)
        
        {\bf parameters:} \\
        \begin{tabular}{| l | l | l | l |}
        \hline
        param  & mode & type & \\
        \hline \hline
        f & I & str or file & file \\ \hline
        \end{tabular}
        
        {\bf Parameter=argument allowed:} \\
        No
        
        {\bf return mode:} \\
        immediate

        {\bf return type:} \\
        log
        
        {\bf Description:} \\
        Return true if the file is a file and not a directory. The argument
        must be immediate mode and may be type "str" of "file". The former is a
        path name string, relative or absolute. The latter is a file variable
        specifying the file.
 


    \section{filelastmod - get the time of last modification}\label{sec:iffilelastmod}
        
        {\bf Call:} \\
        filelastmod(f)
        
        {\bf parameters:} \\
        \begin{tabular}{| l | l | l | l |}
        \hline
        param  & mode & type & \\
        \hline \hline
        f & I & str or file & file \\ \hline
        \end{tabular}
        
        {\bf Parameter=argument allowed:} \\
        No
        
        {\bf return mode:} \\
        immediate

        {\bf return type:} \\
        int
        
        {\bf Description:} \\
        Get the time of last modification of a named file. The time
        is returned as an int representing the time the file was last
        modified, measured in milliseconds since the epoch (00:00:00 GMT,
        January 1, 1970), or 0 if the file does not exist or if an I/O
        error occurs.
        
        The argument must be immediate mode and may be type "str"
        of "file". The former is a path name string, relative or absolute.
        The latter is a file variable specifying the file.
 


    \section{filelength - get the file length in bytes}\label{sec:iffilelength}
        
        {\bf Call:} \\
        filelength(f)
        
        {\bf parameters:} \\
        \begin{tabular}{| l | l | l | l |}
        \hline
        param  & mode & type & \\
        \hline \hline
        f & I & str or file & file \\ \hline
        \end{tabular}
        
        {\bf Parameter=argument allowed:} \\
        No
        
        {\bf return mode:} \\
        immediate

        {\bf return type:} \\
        int
        
        {\bf Description:} \\
        Get the length of a file. The length is the return value. Returns
        0 if the file does not exist.
 
        The argument must be immediate mode and may be type "str"
        of "file". The former is a path name string, relative or absolute.
        The latter is a file variable specifying the file.



    \section{filelist - get a list of files and directories}\label{sec:iffilelist}
        
        {\bf Call:} \\
        filelist(f)
        
        {\bf parameters:} \\
        \begin{tabular}{| l | l | l | l |}
        \hline
        param  & mode & type & \\
        \hline \hline
        f & I & str or file & file \\ \hline
        \end{tabular}
        
        {\bf Parameter=argument allowed:} \\
        No
        
        {\bf return mode:} \\
        immediate

        {\bf return type:} \\
        \verb'[]str'
        
        {\bf Description:} \\
        Get a list of files and directories in a directory.
        
        {\bf f} is the directory and must be immediate mode and may be type "str"
        of "file". The former is a path name string, relative or absolute.
        The latter is a file variable specifying the file.
                
        The names of the files and directories are returned in an array of
        strings. The file and directory names are relative to the specified
        directory, they are not absolute paths.
 


    \section{filelistdirs -  get a list of directories}\label{sec:iffilelistdirs}
        
        {\bf Call:} \\
        filelistdirs(f)
        
        {\bf parameters:} \\
        \begin{tabular}{| l | l | l | l |}
        \hline
        param  & mode & type & \\
        \hline \hline
        f & I & str or file & file \\ \hline
        \end{tabular}
        
        {\bf Parameter=argument allowed:} \\
        No
        
        {\bf return mode:} \\
        immediate

        {\bf return type:} \\
        \verb'[]str'
        
        {\bf Description:} \\
        Get a list of directories in a directory.
        
        {\bf f} is the directory and must be immediate mode and may be type "str"
        of "file". The former is a path name string, relative or absolute.
        The latter is a file variable specifying the file.
                
        The names of the directories are returned in an array of
        strings. The directory names are relative to the specified
        directory, they are not absolute paths.
 


    \section{filelistfiles -  get a list of files}\label{sec:iffilelistfiles}
        
        {\bf Call:} \\
        filelistfiles(f)
        
        {\bf parameters:} \\
        \begin{tabular}{| l | l | l | l |}
        \hline
        param  & mode & type & \\
        \hline \hline
        f & I & str or file & file \\ \hline
        \end{tabular}
        
        {\bf Parameter=argument allowed:} \\
        No
        
        {\bf return mode:} \\
        immediate

        {\bf return type:} \\
        \verb'[]str'
        
        {\bf Description:} \\
        Get a list of files in a directory.
        
        {\bf f} is the directory and must be immediate mode and may be type "str"
        of "file". The former is a path name string, relative or absolute.
        The latter is a file variable specifying the file.
 
        The names of the files are returned in an array of strings. The file
        names are relative to the specified directory, they are not absolute
        paths.


    \section{filemkdir - create a directory}\label{sec:iffilemkdir}
        
        {\bf Call:} \\
        filemkdir(f)
        
        {\bf parameters:} \\
        \begin{tabular}{| l | l | l | l |}
        \hline
        param  & mode & type & \\
        \hline \hline
        f & I & str & file name \\ \hline
        \end{tabular}
        
        {\bf Parameter=argument allowed:} \\
        No
        
        {\bf return mode:} \\
        immediate

        {\bf return type:} \\
        log
        
        {\bf Description:} \\
        Create a directory. Return true if successful.
        {\bf f} is the directory name.
 


    \section{filerename - rename a file or directory}\label{sec:iffilerename}
        
        {\bf Call:} \\
        filerename(fold, fnew)
        
        {\bf parameters:} \\
        \begin{tabular}{| l | l | l | l |}
        \hline
        param  & mode & type & \\
        \hline \hline
        fold & I & str & file name \\ \hline
        fnew & I & str & file name \\ \hline
        \end{tabular}
        
        {\bf Parameter=argument allowed:} \\
        No
        
        {\bf return mode:} \\
        immediate

        {\bf return type:} \\
        log
        
        {\bf Description:} \\
        Rename a file or directory. Return true if successful.
        {\bf fold} is an exisiting file or directory. {\bf fnew}
        is the new name.
 


    \section{filesetreadonly - set a file or directory read-only}\label{sec:iffilesetreadonly}
        
        {\bf Call:} \\
        filesetreadonly(f)
        
        {\bf parameters:} \\
        \begin{tabular}{| l | l | l | l |}
        \hline
        param  & mode & type & \\
        \hline \hline
        f & I & str & file name \\ \hline
        \end{tabular}
        
        {\bf Parameter=argument allowed:} \\
        No
        
        {\bf return mode:} \\
        immediate

        {\bf return type:} \\
        log
        
        {\bf Description:} \\
        set a named file or directory read-only. Return true if successful.
        {\bf f} is the file or directory name.




       

    \section{find - find a substring using a regular expression}\label{sec:iffind}
        
        {\bf Call:} \\
        find(s, re, start)
        
        {\bf parameters:} \\
        \begin{tabular}{| l | l | l | l |}
        \hline
        param  & mode & type & \\
        \hline \hline
        s     & I & str & subject string            \\ \hline
        re    & I & str & regular expression string \\ \hline
        start & I & int & optional starting index   \\ \hline
        \end{tabular}
        
        {\bf Parameter=argument allowed:} \\
        No
        
        {\bf return mode:} \\
        immediate

        {\bf return type:} \\
        \begin{verbatim}
        struct (
            start, "int",
            end,   "int"
        );
        \end{verbatim}
        
        {\bf Description:} \\
        {\bf find()} attempts to match a regular expression string against a
        subject string and returns a struct containing the starting and ending
        indices of the match in the subject string. If there is no match the
        returned indices are zero.
        
        The first argument {\bf s} is the subject string.
        
        The second argument {\bf re} is a Java regular expression - 
        see \autorefph{chap:regexp}.
        
        The third argument {\bf start} is an optional starting index for the
        search and is 0 by default.

        
        See also {\bf match()} \autorefph{sec:ifmatch} for a more complex
        string search.




    \section{first - get the first value of an enum}\label{sec:iffirst}
        
        {\bf Call:} \\
        first(v)
        
        {\bf parameters:} \\
        \begin{tabular}{| l | l | l | l |}
        \hline
        param  & mode & type & \\
        \hline \hline
        v & I V SV S Q & enum or type & enumerated type value \\ \hline
        \end{tabular}
        
        {\bf Parameter=argument allowed:} \\
        No
        
        {\bf return mode:} \\
        immediate

        {\bf return type:} \\
        enum
        
        {\bf Description:} \\
        Returns the first value of an enumerated type. The first value
        is the first in the order the values had as arguments to
        the procedure which created the enum type.
        
        The argument may be either an enum type variable or a variable
        which is an enum type, e.g. for the following variables -
        \begin{lstlisting}[gobble=8, language=threepl]
           enum(e, val1, val2, val3);
           static(s, e);
        \end{lstlisting}
        the argument to {\bf first()} could be either {\bf e} or {\bf s}
        and the result would be the value labelled by identifier {\bf val1}
        of enumerated type {\bf e} in immediate mode.
        
        The argument may be immediate, value, selectvalue, static or queue mode.
        The returned result is immediate mode and the type is the same as the
        argument.


    \section{floatisnegative - sign of target float}\label{sec:iffloatisnegative}
        
        {\bf Call:} \\
        floatisnegative(arg)
        
        {\bf parameters:} \\
        \begin{tabular}{| l | l | l | l |}
        \hline
        param  & mode & type & \\
        \hline \hline
        arg  & V SV S Q IN & float & target float value  \\ \hline
        \end{tabular}
        
        {\bf Parameter=argument allowed:} \\
        No
        
        {\bf return mode:} \\
        value

        {\bf return type:} \\
        log
        
        {\bf Description:} \\
        This function returns the sign bit of a target floating point value
        as type "log".



    \section{floor - round float to nearest integer towards $-\infty$}\label{sec:iffloor}
        
        
        {\bf Call:} \\
        floor(f)
        
        {\bf parameters:} \\
        \begin{tabular}{| l | l | l | l |}
        \hline
        param  & mode & type & \\
        \hline \hline
        f & I & float & \\ \hline
        \end{tabular}
        
        {\bf Parameter=argument allowed:} \\
        No
        
        {\bf return mode:} \\
        immediate

        {\bf return type:} \\
        int
        
        {\bf Description:} \\
        Rounds the floating point argument to the nearest integer towards
        $-\infty$.
        An integer type is returned



    \section{format - format values into a string}\label{sec:ifformat}
        
        
        {\bf Call:} \\
        format(fmt, v1, v2,...)
        
        {\bf parameters:} \\
        \begin{tabular}{| l | l | l | l |}
        \hline
        param  & mode & type & \\
        \hline \hline
        fmt & I & str & format string \\ \hline
        v1  & I &  -  & 1st value to be converted \\ \hline
        v2  & I &  -  & 2nd value to be converted \\ \hline
        .   & I &  -  & ... \\ \hline
        \end{tabular}
        
        {\bf Parameter=argument allowed:} \\
        No
        
        {\bf return mode:} \\
        immediate

        {\bf return type:} \\
        str
        
        {\bf Description:} \\
        Processes a format string consuming the following input parameters as it
        enoucnters conversion specifications. See Java documentation for
        java.lang.String.format .
 

    \section{frequency - get the frequency of a clock variable}\label{sec:iffrequency}
        
        {\bf Call:} \\
        frequency(v)
        
        {\bf parameters:} \\
        \begin{tabular}{| l | l | l | l |}
        \hline
        param  & mode & type & \\
        \hline \hline
        v & C & - & clock variable \\ \hline
        \end{tabular}
        
        {\bf Parameter=argument allowed:} \\
        No
        
        {\bf return mode:} \\
        immediate

        {\bf return type:} \\
        float
        
        {\bf Description:} \\
        Returns the frequency of a clock variable in Hz. If the frequency
        is variable this is the maximum frequency.
        
        See also inbuilt functions - \\
        {\bf period()} \autorefph{sec:ifperiod} \\
        {\bf dutycycle()} \autorefph{sec:ifdutycycle}


    \section{get - get a list or map entry}\label{sec:ifget}
        
        {\bf Call:} \\
        get(l, index) \\
        get(m, key) \\
        get(m, index)
        
        {\bf parameters:} \\
        \begin{tabular}{| l | l | l | l |}
        \hline
        param  & mode & type & \\
        \hline \hline
        l     & I & list & a list \\ \hline
        m     & I & map  & a map \\ \hline
        index & I & int  & a list or map index \\ \hline
        key   & I & str  & a map key \\ \hline
        \end{tabular}
        
        {\bf Parameter=argument allowed:} \\
        No
        
        {\bf return mode:} \\
        value

        {\bf return type:} \\
        type of list or map entry
        
        {\bf Description:} \\
        This function can be used to get a list or a map entry.
        
        For a list {\bf get(l, index)} returns the item at the position in
        the list given by the index. The index must be in the range 0 to
        entries(l) - 1 inclusive.
        
        For a map {\bf get(m, key)} returns the item matching the string key.
        If there is no such entry, null is returned.
        The items are always stored in the lexicographic order of the keys. The
        call {\bf get(m, index)} returns the key at the position in the ordered
        keys given by the index. The index must be in the range 0 to entries(l)
        - 1 inclusive. If the index is out of range a fatal error occurs. To get
        the value instead of the key at an index use {\bf get(get(m, index))}.



%    \section{getfield - get the nth field from a struct}\label{sec:ifgetfield}
%        
%        {\bf Call:} \\
%        getfield(v, n)
%        
%        {\bf parameters:} \\
%        \begin{tabular}{| l | l | l | l |}
%        \hline
%        param  & mode & type & \\
%        \hline \hline
%        v & I V SV S Q & struct & variable  \\ \hline
%        n & I          & int    & field     \\ \hline
%        \end{tabular}
%        
%        {\bf Parameter=argument allowed:} \\
%        No
%        
%        {\bf return mode:} \\
%        mode of the struct
%
%        {\bf return type:} \\
%        type of the field
%        
%        {\bf Description:} \\
%        This function returns the value of the field specified by the second
%        argument from the struct given as the first argument.


    \section{hadeof - had EOF on the last file read}\label{sec:ifhadeof}
        
        {\bf Call:} \\
        hadeof(file)
        
        {\bf parameters:} \\
        \begin{tabular}{| l | l | l | l |}
        \hline
        param  & mode & type & \\
        \hline \hline
        file & I & file & open input file  \\ \hline
        \end{tabular}
        
        {\bf Parameter=argument allowed:} \\
        No
        
        {\bf return mode:} \\
        immediate

        {\bf return type:} \\
        log
        
        {\bf Description:} \\
        This function returns true if the most recent read was at
        end-of-file.


%    \section{hasfield - check if a struct has a specific field}\label{sec:ifhasfield}
%        
%        {\bf Call:} \\
%        hasfield(v, f)
%        
%        {\bf parameters:} \\
%        \begin{tabular}{| l | l | l | l |}
%        \hline
%        param  & mode & type & \\
%        \hline \hline
%        v & I V SV S Q & - & variable  \\ \hline
%        f & -          & - & field     \\ \hline
%        \end{tabular}
%        
%        {\bf Parameter=argument allowed:} \\
%        No
%        
%        {\bf return mode:} \\
%        immediate
%
%        {\bf return type:} \\
%        log
%        
%        {\bf Description:} \\
%        This function returns true if the field identifier given by the second
%        argument is included in the struct given as the first argument.
%        The first argument may be a struct variable (which is a variable of type
%        {\bf type}) or it may be an immediate or target variable that is a
%        struct type. For example in the code -
%        \begin{lstlisting}[gobble=8, language=threepl]
%           struct(s,
%               f1, "log",
%               f2, "int"
%           );
%           var(v, s);
%           ...hasfield(s, f1).....
%           ...hasfield(v, f1).....
%        \end{lstlisting}
%        the first argument to {\bf hasfield()} can be either the struct
%        variable itself or the variable that is the struct type.
%        
%        If the first argument is not an appropriate type {\bf hasfield()}
%        will simply return false.
 

    \section{hasnext - is there a next value of an enum}\label{sec:ifhasnext}
        
        {\bf Call:} \\
        hasnext(v)
        
        {\bf parameters:} \\
        \begin{tabular}{| l | l | l | l |}
        \hline
        param  & mode & type & \\
        \hline \hline
        v & I V SV S Q & enum & enumerated type value \\ \hline
        \end{tabular}
        
        {\bf Parameter=argument allowed:} \\
        No
        
        {\bf return mode:} \\
        immediate or value

        {\bf return type:} \\
        log
        
        {\bf Description:} \\
        Returns true if there is a next value of an enumerated type. The next
        value is the successor in the order the values had as arguments to
        the procedure which created the enum type. An alternative view is
        that the function returns false if the argument value is the last
        value in the ordering of the enumerated type, otherwise it returns true.
        
        If the argument is immediate mode the returned result is immediate
        mode and the type is log.
        
        If the argument is a target mode (value, selectvalue, static or queue)
        the returned result is value mode and the type log. For an argument
        which is a queue the result will only be available when the queue is
        available.
 


    \section{hasqueuereads - does an argument contain queue reads}\label{sec:ifhasqueuereads}
        
        {\bf Call:} \\
        hasqueuereads(arg)
        
        {\bf parameters:} \\
        \begin{tabular}{| l | l | l | l |}
        \hline
        param  & mode & type & \\
        \hline \hline
        arg & - & any & expression  \\ \hline
        \end{tabular}
        
        {\bf Parameter=argument allowed:} \\
        No
        
        {\bf return mode:} \\
        immediate

        {\bf return type:} \\
        log
        
        {\bf Description:} \\
        This function returns true if the argument is a target variable or
        expression containing queue reads (not queue examines). If the argument
        only contains constants or immediate variables it returns false. If the
        argument is a clock, cmemory or rmemory mode variable the function will
        return false.
        
        This function is used to determine if a value argument to a module,
        procedure or function contains any queue reads to avoid further
        calls to inbuilt or library modules, procedures or functions
        which prohibit queues in their arguments.
        
        If called with no arguments this function checks all input parameters
        to the current module, procedure or function and returns true if
        any are queues or contain queue reads.
 


    \section{hasprev - is there a previous value of an enum}\label{sec:ifhasprev}
        
        {\bf Call:} \\
        hasprev(v)
        
        {\bf parameters:} \\
        \begin{tabular}{| l | l | l | l |}
        \hline
        param  & mode & type & \\
        \hline \hline
        v & I V SV S Q & enum & enumerated type value \\ \hline
        \end{tabular}
        
        {\bf Parameter=argument allowed:} \\
        No
        
        {\bf return mode:} \\
        immediate or value

        {\bf return type:} \\
        log
        
        {\bf Description:} \\
        Returns true if there is a previous value of an enumerated type. The
        previous value is the predecessor in the order the values had as
        arguments to the procedure which created the enum type. An alternative
        view is that the function returns false if the argument value is the
        first value in the ordering of the enumerated type, otherwise it
        returns true.
        
        If the argument is immediate mode the returned result is immediate
        mode and the type is log.
        
        If the argument is a target mode (value, selectvalue, static or queue)
        the returned result is value mode and the type log. For an argument
        which is a queue the result will only be available when the queue is
        available.

       

    \section{head - get the head of a string}\label{sec:ifhead}
        
        {\bf Call:} \\
        head(s1, i2)
        
        {\bf parameters:} \\
        \begin{tabular}{| l | l | l | l |}
        \hline
        param  & mode & type & \\
        \hline \hline
        s1 & I & str & subject string   \\ \hline
        i2 & I & int & length of head   \\ \hline
        \end{tabular}
        
        {\bf Parameter=argument allowed:} \\
        No
        
        {\bf return mode:} \\
        immediate

        {\bf return type:} \\
        str
        
        {\bf Description:} \\
        Returns the head of the string first argument, the length of the
        head string being given by the integer second argument. If the
        length requested is greater than the string length a fatal error
        occurs.


    \section{identifier - get the identifier of a variable}\label{sec:ifidentifier}
        
        {\bf Call:} \\
        identifier(ident, o)
        
        {\bf parameters:} \\
        \begin{tabular}{| l | l | l | l |}
        \hline
        param  & mode & type & \\
        \hline \hline
        ident & - &  -  & identifier or identifier function \\ \hline
        o     & I & str & optional type                     \\ \hline
        \end{tabular}
        
        {\bf Parameter=argument allowed:} \\
        No
        
        {\bf return mode:} \\
        immediate

        {\bf return type:} \\
        str
        
        {\bf Description:} \\
        Return a string which is the identifier of the variable. This has 1 or
        2 arguments. The 1st or only argument is the variable. The optional 2nd
        argument selects the form of the returned identifier string.
        
        The argument may also be one of the functions which returns a variable
        identifier such as {\bf getclock()}, {\bf currentclock()},
        {\bf prevclock()}, {\bf domain}, {\bf readdomain} and {\bf writedomain}.
        A null return from one of these functions will result in the string
        "NULL".
        
        {\bf o} selects the form in which the identifier string is returned as follows
        
        {\bf o} omitted - the literal identifier string is returned
        
        {\bf o} = "sliteral" - the identifier is prepended by the identifier of the module,
        procedure or function in which it has been created, except where the variable
        was created within a nested code block.In the case of file variables the
        display format is {\em file:var} and for global variables {\em glob:var}.
        
        {\bf o} = "ext" - the complete unambiguous variable name is returned.
        
        {\bf o} = "real" - for a parameter the argument variable identifier will be returned.
        For an assigned clock variable the identifier of the RHS variable of the assignment
        will be returned. Otherwise the variable identifier string is returned.
        
        {\bf o} = "chain" - as for "real" but both the parameter identifier and the argument identifier
        will be returned separated by \verb'->'. The identifiers are preceded by the name of the
        module, procedure or function in which they were created.


%    \section{imtotargtype - convert an immediate type to a target type}\label{sec:ifimtotargtype}
%        
%        {\bf Call:} \\
%        imtotargtype(expr)
%        
%        {\bf parameters:} \\
%        \begin{tabular}{| l | l | l | l |}
%        \hline
%        param  & mode & type & \\
%        \hline \hline
%        expr & any & any & value \\ \hline
%        \end{tabular}
%        
%        {\bf Parameter=argument allowed:} \\
%        No
%        
%        {\bf return mode:} \\
%        immediate
%
%        {\bf return type:} \\
%        str
%        
%        {\bf Description:} \\
%        Return the type of an immediate variable as a target type. The type is a
%        string in the same format as is used for type declarations. The type
%        will be converted to target type using widths derived from the immediate
%        variable value(s).
%        
%        {\bf THIS FUNCTION HAS NOT YET BEEN COMPLETED AND IS A TEMPLATE
%        WHICH AT PRESENT WILL GIVE A FATAL ERROR!}
%        
%        See also - \\
%        {\bf type()} \autorefph{sec:iftype} \\
%        {\bf arraytype()} \autorefph{sec:ifarraytype} \\
%        {\bf typestring()} \autorefph{sec:iftypestring} \\
%        {\bf primtypestring()} \autorefph{sec:ifprimtypestring} \\
%        {\bf targtoimtype()} \autorefph{sec:iftargtoimtype} \\
%        {\bf typeval()} \autorefph{sec:iftypeval} \\
%        {\bf typevalstring()} \autorefph{sec:iftypevalstr} \\
%        {\bf typeisequal()} \autorefph{sec:iftypeiseq} \\






    \section{includedirs - get the directories searched for module and source files}\label{sec:ifincldirs}
        
        {\bf Call:} \\
        includedirs()
        
        {\bf parameters:} \\
        None
        
        {\bf Parameter=argument allowed:} \\
        No
        
        {\bf return mode:} \\
        immediate

        {\bf return type:} \\
        \verb'[]str'
        
        {\bf Description:} \\
        Return a string array containing the absolute paths of directories
        searched for module and source included files.





    \section{inttobinstring - get a binary string representing an integer}\label{sec:ifinttobinstring}
        
        {\bf Call:} \\
        inttohexstring(v, l)
        
        {\bf parameters:} \\
        \begin{tabular}{| l | l | l | l |}
        \hline
        param  & mode & type & \\
        \hline \hline
        v & I & bits, int & value           \\ \hline
        l & I & int       & minimum length  \\ \hline
        \end{tabular}
        
        {\bf Parameter=argument allowed:} \\
        No
        
        {\bf return mode:} \\
        immediate

        {\bf return type:} \\
        str
        
        {\bf Description:} \\
        Return a binary string which represents the value of the first
        argument. The first argument must be type bits or int. The optional
        second argument gives the minimum length of the string, the result being
        zero-padded on the left if necessary.





    \section{inttohexstring - get a hex string representing an integer}\label{sec:ifinttohexstring}
        
        {\bf Call:} \\
        inttohexstring(v, l)
        
        {\bf parameters:} \\
        \begin{tabular}{| l | l | l | l |}
        \hline
        param  & mode & type & \\
        \hline \hline
        v & I & bits, int & value           \\ \hline
        l & I & int       & minimum length  \\ \hline
        \end{tabular}
        
        {\bf Parameter=argument allowed:} \\
        No
        
        {\bf return mode:} \\
        immediate

        {\bf return type:} \\
        str
        
        {\bf Description:} \\
        Return a hexadecimal string which represents the value of the first
        argument. The first argument must be type bits or int. The optional
        second argument gives the minimum length of the string, the result being
        zero-padded on the left if necessary.
       

    \section{isnull - determine if null}\label{sec:ifisnull}
        
        {\bf Call:} \\
        isnull(v)
        
        {\bf parameters:} \\
        \begin{tabular}{| l | l | l | l |}
        \hline
        param  & mode & type & \\
        \hline \hline
        v & V & any & variable or function \\ \hline
        \end{tabular}
        
        {\bf Parameter=argument allowed:} \\
        No
        
        {\bf return mode:} \\
        immediate

        {\bf return type:} \\
        log
        
        {\bf Description:} \\
        Returns true if the argument is null. This function provides a
        null test for values which might not be primitive. The use of
        operators {\bf ==} and {\bf !=} causes a fatal error if one of
        the operands is not a primitive type.
       

    \section{isread - determine when a queue variable is read}\label{sec:ifisread}
        
        {\bf Call:} \\
        isread(v)
        
        {\bf parameters:} \\
        \begin{tabular}{| l | l | l | l |}
        \hline
        param  & mode & type & \\
        \hline \hline
        v & Q & any & queue mode variable \\ \hline
        \end{tabular}
        
        {\bf Parameter=argument allowed:} \\
        No
        
        {\bf return mode:} \\
        value

        {\bf return type:} \\
        log
        
        {\bf Description:} \\
        Returns a target value of type "log" which is true on the clock cycle in
        which a read from a queue variable occurs. The argument must be a queue variable.
        
        The queue variable argument must not have subscripts or fields as a read
        from a queue variable always reads all components of that variable so
        it would be meaningless to specify some components of the queue variable.
        
        The argument variable read domain must be the same as the clock domain
        of the call to {\bf isread()}. If a queue mode variable does
        not yet have a (read) clock domain established then the call to
        {\bf isread()} will establish that domain.
        
        See also - \\
        {\bf iswritten()} \autorefph{sec:ifiswritten} \\
        {\bf waswritten()} \autorefph{sec:ifwaswritten}


    \section{issigned - determine if a type is signed}\label{sec:ifissigned}
        
        {\bf Call:} \\
        issigned(arg)
        
        {\bf parameters:} \\
        \begin{tabular}{| l | l | l | l |}
        \hline
        param  & mode & type & \\
        \hline \hline
        arg  & V SV S Q IN I & any &   \\ \hline
        \end{tabular}
        
        {\bf Parameter=argument allowed:} \\
        No
        
        {\bf return mode:} \\
        immediate

        {\bf return type:} \\
        log
        
        {\bf Description:} \\
        This function returns true if the type of the argument is
        "INT:" or "FIXED:". If the argument is immediate mode its type
        must be "type" in which case its type value will be tested
        for "INT:" or "FIXED:".


    \section{isstruct - determine if a variable type is a struct}\label{sec:ifisstruct}
        
        {\bf Call:} \\
        isstruct(v)
        
        {\bf parameters:} \\
        \begin{tabular}{| l | l | l | l |}
        \hline
        param  & mode & type & \\
        \hline \hline
        v & I V SV S Q & struct & variable  \\ \hline
        \end{tabular}
        
        {\bf Parameter=argument allowed:} \\
        No
        
        {\bf return mode:} \\
        immediate

        {\bf return type:} \\
        log
        
        {\bf Description:} \\
        This function returns true if the type of the argument is a struct.
       

    \section{istargconst - determine if a value mode variable is a constant}\label{sec:ifistargconst}
        
        {\bf Call:} \\
        istargconst(v)
        
        {\bf parameters:} \\
        \begin{tabular}{| l | l | l | l |}
        \hline
        param  & mode & type & \\
        \hline \hline
        v & V & prim & value mode variable \\ \hline
        \end{tabular}
        
        {\bf Parameter=argument allowed:} \\
        No
        
        {\bf return mode:} \\
        immediate

        {\bf return type:} \\
        log
        
        {\bf Description:} \\
        Returns true if the value mode variable {\bf v} is a target constant.

        

    \section{iswritten - determine when a static or queue variable is written}\label{sec:ifiswritten}
        
        {\bf Call:} \\
        iswritten(v)
        
        {\bf parameters:} \\
        \begin{tabular}{| l | l | l | l |}
        \hline
        param  & mode & type & \\
        \hline \hline
        v & S Q & any & static or queue mode variable \\ \hline
        \end{tabular}
        
        {\bf Parameter=argument allowed:} \\
        No
        
        {\bf return mode:} \\
        value

        {\bf return type:} \\
        log
        
        {\bf Description:} \\
        Returns a target value of type "log" which is true on the clock cycle in
        which a write to a static or queue variable or a part of a static
        variable occurs. The argument must be a static or queue variable.

        A static variable argument may have struct fields or array subscripts in
        which case the returned value will be true following a write to any
        field or array member within those specified by the argument.

        A queue variable argument must not have subscripts or fields as a write
        to a queue variable always writes all components of that variable so it
        would be meaningless to specify some components of the queue variable.

        The argument variable write domain must be the same as the clock domain
        of the call to {\bf iswritten()}. If a static or queue mode variable does
        not yet have a (write) clock domain established then the call to {\bf
        iswritten()} will establish that domain.

        {\bf This function should be used with caution. It must not be used in a
        conditional assignment to the variable which is the argument to the
        function as this is a logical contradiction and would lead to a closed
        loop in the logic!}

        See also inbuilt function \\
        {\bf waswritten()}
        \autorefph{sec:ifwaswritten} to signal on the cycle following the one on
        which the write occurs. \\
        {\bf isread()} \autorefph{sec:ifisread} can be
        used to signal on the cycle a queue is read.


    \section{last - get the last value of an enum}\label{sec:iflast}
        
        {\bf Call:} \\
        last(v)
        
        {\bf parameters:} \\
        \begin{tabular}{| l | l | l | l |}
        \hline
        param  & mode & type & \\
        \hline \hline
        v & I V SV S Q & enum or type & enumerated type value \\ \hline
        \end{tabular}
        
        {\bf Parameter=argument allowed:} \\
        No
        
        {\bf return mode:} \\
        immediate

        {\bf return type:} \\
        enum
        
        {\bf Description:} \\
        Returns the last value of an enumerated type. The last value
        is the last in the order the values had as arguments to
        the procedure which created the enum type.
        
        The argument may be either an enum type variable or a variable
        which is that enum type, e.g. for the following variables -
        \begin{lstlisting}[gobble=8, language=threepl]
           enum(e, val1, val2, val3);
           static(s, e);
        \end{lstlisting}
        the argument to {\bf first()} could be either {\bf e} or {\bf s}
        and the result would be the value labelled by identifier {\bf val3}
        of enumerated type {\bf e} in immediate mode.
        
        The argument may be immediate, value, selectvalue, static or queue mode.
        The returned result is immediate mode and the type is the same as the
        argument.
 

    \section{listattributes - get a list of variable attributes}\label{sec:iflistattrs}
        
        {\bf Call:} \\
        listattributes(v)
        
        {\bf parameters:} \\
        \begin{tabular}{| l | l | l | l |}
        \hline
        param  & mode & type & \\
        \hline \hline
        v      & any  & any  & variable identifier \\ \hline
        \end{tabular}
        
        {\bf Parameter=argument allowed:} \\
        No
        
        {\bf return mode:} \\
        immediate

        {\bf return type:} \\
        \verb'[n](key, str, mode, str, value, str)'
        
        {\bf Description:} \\
        Returns a list of attributes for a value. The returned value is an
        array, each member being a struct. The struct has three fields, {\em
        key} is a string giving the attribute name, {\em mode} is a string
        giving the mode of the attribute value and {\em value} is a string
        representation of the value or the identifier in the case of a clock
        domain.
        
        The argument is the variable identifier.
        
        The struct returned by the function can be used without requiring the
        structure definition by the following -
        \begin{lstlisting}[gobble=8, language=threepl]
           var(ld, "empty", listattributes());
        \end{lstlisting}
        after which {\bf ld[i].key}, {\bf ld[i].mode} etc can be used.
        
        See also library procedure {\bf printattributes()} \autorefph{sec:lpprintattrs}.


    \section{listvars - get a list of variables}\label{sec:iflistvars}
        
        {\bf Call:} \\
        listvars(scope, allsegs)
        
        {\bf parameters:} \\
        \begin{tabular}{| l | l | l | l |}
        \hline
        param  & mode & type & \\
        \hline \hline
        scope & I & str & optional scope \\ \hline
        allsegs & I & str & optional flag \\ \hline
        \end{tabular}
        
        {\bf Parameter=argument allowed:} \\
        No
        
        {\bf return mode:} \\
        immediate

        {\bf return type:} \\
        \verb'[n](scope, str, ident, str, var, ->)'
        
        {\bf Description:} \\
        Returns a list of variables, their scope and type. The returned value is
        an array, each member being a struct. The struct has three fields, {\em
        scope} is a string giving the scope, {\em ident} is a string giving the
        identifier and {\em var} is a pointer to the variable.

        The optional first argument, {\bf scope},  is a string limiting a scope.
        By default all block and local variables, the file scope variables and
        global scope variables will be listed.
        
        Argument "global" causes only global variables to be listed.

        Argument "file" causes only file scope variables to be listed.

        Argument "files" causes variables in all file scopes to be listed.

        Argument "local" causes only local variables to be listed, i.e. only
        those declared within the current module, procedure or function, and
        not variables in nested control structures.

        Argument "calling" causes only variables in the calling scope to be
        listed.
        
        Argument "" selects the default, i.e. all block and local variables, the
        file scope variables and global scope variables will be listed.
        
        The optional second argument, {\bf allsegs}, controls which file scopes
        are examined. If true, variables in all file scopes are listed otherwise
        only variables in the current file scope are listed.
        
        The variables within each scope are listed in lexicographic order.
        
        The struct returned by the function can be used without requiring the
        structure definition by the following -
        \begin{lstlisting}[gobble=8, language=threepl]
           var(lv, "empty", listvars());
        \end{lstlisting}
        after which {\bf lv[i].scope}, {\bf lv[i].ident} etc can be used.
        
        See also library procedure {\bf printvars()} \autorefph{sec:lpprintvars}.
 

    \section{location - get the current source location}\label{sec:iflocation}
        
        {\bf Call:} \\
        location()
        
        {\bf parameters:} \\
        none
        
        {\bf return mode:} \\
        immediate

        {\bf return type:} \\
        str
        
        {\bf Description:} \\
        Returns the current location in the source code as a string. The
        location gives the source file name, line number and column number.
        
        See also inbuilt procedure {\bf trace()} \autorefph{sec:iptrace}
        which prints the location and stack trace to standard out.


    \section{log - natural logarithm}\label{sec:iflog}
            
        {\bf Call:} \\
        log(f)
        
        {\bf parameters:} \\
        \begin{tabular}{| l | l | l | l |}
        \hline
        param  & mode & type & \\
        \hline \hline
        f & I & float & \\ \hline
        \end{tabular}
        
        {\bf Parameter=argument allowed:} \\
        No
        
        {\bf return mode:} \\
        immediate

        {\bf return type:} \\
        float
        
        {\bf Description:} \\
        Returns the natural logarithm (base {\em e}) of the immediate
        argument.


    \section{mantissa - target float mantissa}\label{sec:ifmantissa}
        
        {\bf Call:} \\
        mantissa(arg)
        
        {\bf parameters:} \\
        \begin{tabular}{| l | l | l | l |}
        \hline
        param  & mode & type & \\
        \hline \hline
        arg  & V SV S Q IN & float & target float value  \\ \hline
        \end{tabular}
        
        {\bf Parameter=argument allowed:} \\
        No
        
        {\bf return mode:} \\
        value

        {\bf return type:} \\
        uint
        
        {\bf Description:} \\
        This function returns the mantissa of a target floating point
        value. The mantissa is return as type uint: whose width is
        equal to the width of the mantissa.


    \section{mantissawidth - width of target float mantissa}\label{sec:ifmantissawidth}
        
        {\bf Call:} \\
        mantissawidth(arg)
        
        {\bf parameters:} \\
        \begin{tabular}{| l | l | l | l |}
        \hline
        param  & mode & type & \\
        \hline \hline
        arg  & I V SV S Q IN & type float & target float type or value  \\ \hline
        \end{tabular}
        
        {\bf Parameter=argument allowed:} \\
        No
        
        {\bf return mode:} \\
        immediate

        {\bf return type:} \\
        uint
        
        {\bf Description:} \\
        This function returns the number of bits in the mantissa of
        a target floating point value. The argument may be a target floating
        point value or it may be a target floating point type.




       

    \section{match - find substrings using a regular expression}\label{sec:ifmatch}
        
        {\bf Call:} \\
        match(s, re, start)
        
        {\bf parameters:} \\
        \begin{tabular}{| l | l | l | l |}
        \hline
        param  & mode & type & \\
        \hline \hline
        s     & I & str & optional subject string               \\ \hline
        re    & I & str & optional regular expression string    \\ \hline
        start & I & int & optional starting index               \\ \hline
        \end{tabular}
        
        {\bf Parameter=argument allowed:} \\
        No
        
        {\bf return mode:} \\
        immediate

        {\bf return type:} \\
        "list"
        
        {\bf Description:} \\
        {\bf match()} attempts to match a regular expression string against a
        subject string and returns a list of the groups of matched substrings.
        If there is no match the returned list is empty.
        
        The first argument {\bf s} is the subject string.
        
        The second argument {\bf re} is a Java regular expression - 
        see \autorefph{chap:regexp}.
        
        The third argument {\bf start} is an optional starting index for the
        search and is 0 by default.
        
        To perform further matches on the same subject with the same pattern
        from the end of the previous match the function is called again with
        but with no arguments.
        
        Group 0 in the returned list is the entire substring matched. Following
        groups are substrings matched by parenthesised expressions in the
        regular expression.
        
        If a match fails an empty list is returned.
        
        See also {\bf find()} \autorefph{sec:iffind} for a simpler
        string search.




    \section{matched - determine if a parameter has an argument}\label{sec:ifmatched}
        
        
        {\bf Call:} \\
        matched(ident)
        
        {\bf parameters:} \\
        \begin{tabular}{| l | l | l | l |}
        \hline
        param  & mode & type & \\
        \hline \hline
        ident & - &  -  & identifier \\ \hline
        \end{tabular}
        
        {\bf Parameter=argument allowed:} \\
        No
        
        {\bf return mode:} \\
        immediate

        {\bf return type:} \\
        log
        
        {\bf Description:} \\
        Returns true if the identifier is that of a module, procedure or
        function parameter which has been passed an argument. Also returns true
        if the identifier is that of a variable which is not a module, procedure
        or function parameter.

        Returns false if the identifier is that of a parameter which has not
        been passed an argument.
        
        Note that a parameter which is not passed an argument becomes a local
        variable. If that parameter is passed as an argument to a further
        call there is no indication in the deeper scope that the argument
        two levels back was omitted. To resolve this problem inbuilt function
        {\bf nullarg()} \autorefph{sec:ifnullarg} can be used.
        
        A fatal error occurs if the identifier is not recognised.


    \section{max - maximum of integers}\label{sec:ifmax}
        
        
        {\bf Call:} \\
        max(i1, i2, ...)
        
        {\bf parameters:} \\
        \begin{tabular}{| l | l | l | l |}
        \hline
        param  & mode & type & \\
        \hline \hline
        i1 & I & int & integer value    \\ \hline
        i2 & I & int & integer value    \\ \hline
        .  & I & int & etc              \\ \hline
        \end{tabular}
        
        {\bf Parameter=argument allowed:} \\
        No
        
        {\bf return mode:} \\
        immediate

        {\bf return type:} \\
        int
        
        {\bf Description:} \\
        Returns the maximum of the integer arguments. The arguments may
        be negative.


    \section{maxvalue - maximum value of a target type}\label{sec:ifmaxvalue}
        
        
        {\bf Call:} \\
        maxvalue(v)
        
        {\bf parameters:} \\
        \begin{tabular}{| l | l | l | l |}
        \hline
        param  & mode & type & \\
        \hline \hline
        v & target & uint, int, enum & type    \\ \hline
        \end{tabular}
        
        {\bf Parameter=argument allowed:} \\
        No
        
        {\bf return mode:} \\
        value

        {\bf return type:} \\
        type specified by argument
        
        {\bf Description:} \\
        Returns the maximum value of the variable argument. The argument must
        specify a primitive target type. It may be a structured type with fields
        and/or subscripts appended such that it represent a primitive type. The
        argument mode cannot be CLOCK, CMEMORY or RMEMORY mode. The maximum type
        size is 63 bits.

        The argument may be a target-mode variable in which case the maximum
        value the variable can represent will be returned.

        The variable may also be an immediate modeof type "type" in which case
        the maximum value of the type will be returned. A type variable may also
        have fields and/or subscripts where the type is structured.

        The argument may also be an immediate variable, expression or
        literalconstant of type "str" in which case it will be interpreted as a
        type.

        The types currently allowed are "int:", "uint:", "fixed", "ufixed",
        "float" and "enum". In the latter case the largest ordinal is returned.
        
        See also {\bf minvalue()} \autorefph{sec:ifminvalue} for the minimum
        value.



    \section{memoryaddresstype - get the address type of a memory variable}\label{sec:ifmemoryaddresstype}
        
        {\bf Call:} \\
        memoryaddresstype(v)
        
        {\bf parameters:} \\
        \begin{tabular}{| l | l | l | l |}
        \hline
        param  & mode & type & \\
        \hline \hline
        v & M & - & cmemory or rmemory mode variable \\ \hline
        \end{tabular}
        
        {\bf Parameter=argument allowed:} \\
        No
        
        {\bf return mode:} \\
        immediate

        {\bf return type:} \\
        type
        
        {\bf Description:} \\
        Returns the address type of a cmemory or rmemory mode variable.

        See also inbuilt function
        {\bf memorydatatype()} \autorefph{sec:ifmemorydatatype}.

    \section{memorydatatype - get the data type of a memory variable}\label{sec:ifmemorydatatype}
        
        {\bf Call:} \\
        memorydatatype(v)
        
        {\bf parameters:} \\
        \begin{tabular}{| l | l | l | l |}
        \hline
        param  & mode & type & \\
        \hline \hline
        v & M & - & cmemory or rmemory mode variable \\ \hline
        \end{tabular}
        
        {\bf Parameter=argument allowed:} \\
        No
        
        {\bf return mode:} \\
        immediate

        {\bf return type:} \\
        type
        
        {\bf Description:} \\
        Returns the data type of a cmemory or rmemory mode variable.

        See also inbuilt function
        {\bf memoryaddresstype()} \autorefph{sec:ifmemoryaddresstype}.



    \section{min - minimum of integers}\label{sec:ifmin}
        
        
        {\bf Call:} \\
        min(i1, i2, ...)
        
        {\bf parameters:} \\
        \begin{tabular}{| l | l | l | l |}
        \hline
        param  & mode & type & \\
        \hline \hline
        i1 & I & int & integer value    \\ \hline
        i2 & I & int & integer value    \\ \hline
        .  & I & int & etc              \\ \hline
        \end{tabular}
        
        {\bf Parameter=argument allowed:} \\
        No
        
        {\bf return mode:} \\
        immediate

        {\bf return type:} \\
        int
        
        {\bf Description:} \\
        Returns the minimum of the integer arguments. The arguments may
        be negative.


    \section{minvalue - minimum value of a target type}\label{sec:ifminvalue}
        
        
        {\bf Call:} \\
        minvalue(v)
        
        {\bf parameters:} \\
        \begin{tabular}{| l | l | l | l |}
        \hline
        param  & mode & type & \\
        \hline \hline
        v & target & uint, int, enum & type    \\ \hline
        \end{tabular}
        
        {\bf Parameter=argument allowed:} \\
        No
        
        {\bf return mode:} \\
        value

        {\bf return type:} \\
        type specified by argument
        
        {\bf Description:} \\
        Returns the minimum value of the variable argument. The argument must
        specify a primitive target type. It may be a structured type with fields
        and/or subscripts appended such that it represent a primitive type. The
        argument mode cannot be CLOCK, CMEMORY or RMEMORY mode. The maximum type
        size is 63 bits.

        The argument may be a target-mode variable in which case the minimum
        value the variable can represent will be returned.

        The variable may also be an immediate modeof type "type" in which case
        the minimum value of the type will be returned. A type variable may also
        have fields and/or subscripts appended where the type is structured to
        reduce to a primitive type.

        The argument may also be an immediate variable, expression or literal
        constant of type "str" in which case it will be interpreted as a type.

        The types currently allowed are "int:", "uint:", "fixed", "ufixed",
        "float" and "enum". In the latter case the largest ordinal is returned.
        
        See also {\bf maxvalue()} \autorefph{sec:ifmaxvalue} for the maximum
        value.


    \section{mod - modulus of integer or float}\label{sec:ifmod}
        
        
        {\bf Call:} \\
        mod(v)
        
        {\bf parameters:} \\
        \begin{tabular}{| l | l | l | l |}
        \hline
        param  & mode & type & \\
        \hline \hline
        v & I & int, float &  \\ \hline
        \end{tabular}
        
        {\bf Parameter=argument allowed:} \\
        No
        
        {\bf return mode:} \\
        immediate

        {\bf return type:} \\
        int or float
        
        {\bf Description:} \\
        Returns the modulus of the integer or floating point argument.
        The result is integer or floating to match the argument type.


    \section{mode - the mode of a variable}\label{sec:ifmode}
        
        {\bf Call:} \\
        mode(ident)
        
        {\bf parameters:} \\
        \begin{tabular}{| l | l | l | l |}
        \hline
        param  & mode & type & \\
        \hline \hline
        ident & - &  -  & variable identifier \\ \hline
        \end{tabular}
        
        {\bf Parameter=argument allowed:} \\
        No
        
        {\bf return mode:} \\
        immediate

        {\bf return type:} \\
        str
        
        {\bf Description:} \\
        Return the mode of a variable as an string. Values are -\\
        \begin{tabular}{| l |}
        \hline
        "immediate"     \\ \hline
        "value"         \\ \hline
        "selectvalue"   \\ \hline
        "static"        \\ \hline
        "queue"          \\ \hline
        "priority"      \\ \hline
        "input"         \\ \hline
        "output"        \\ \hline
        "cmemory"       \\ \hline
        "rmemory"       \\ \hline
        "clock"         \\ \hline
        \end{tabular}

    \section{next - get the next value of an enum}\label{sec:ifnext}
        
        {\bf Call:} \\
        next(v)
        
        {\bf parameters:} \\
        \begin{tabular}{| l | l | l | l |}
        \hline
        param  & mode & type & \\
        \hline \hline
        v & I V SV S Q & enum & enumerated type value \\ \hline
        \end{tabular}
        
        {\bf Parameter=argument allowed:} \\
        No
        
        {\bf return mode:} \\
        immediate or value

        {\bf return type:} \\
        enum
        
        {\bf Description:} \\
        Returns the next value of an enumerated type. The next value
        is the successor in the order the values had as arguments to
        the procedure which created the enum type.
        
        If the argument is immediate mode the returned result is immediate
        mode and the type is the same as the argument. If there is no next
        value a fatal error occurs.
        
        If the argument is a target mode (value, selectvalue, static or queue)
        the returned result is value mode and the type is the same as the
        argument. For an argument which is a queue the result will only be
        available when the queue is available. If there is no next
        value a value with an ordinal of 0 is returned.


    \section{nullarg - omit an unmatched parameter}\label{sec:ifnullarg}
        
        
        {\bf Call:} \\
        nullarg(ident)
        
        {\bf parameters:} \\
        \begin{tabular}{| l | l | l | l |}
        \hline
        param  & mode & type & \\
        \hline \hline
        ident & - &  -  & identifier \\ \hline
        \end{tabular}
        
        {\bf Parameter=argument allowed:} \\
        No
        
        {\bf return mode:} \\
        any

        {\bf return type:} \\
        any
        
        {\bf Description:} \\
        This function is used to wrap an argument to a module, procedure or function
        where the argument is itself a module, procedure or function parameter.
        The argument to nullarg() must be a parameter. It may have subscripts,
        fields or pre or post increment or decrement operators.

        If the argument to nullarg() is a parameter which has been matched by an
        argument then the argument is passed in place of the call to nullarg().
        
        If the argument to nullarg() is a parameter which has not been matched by
        an argument then no argument is passed in place of the call to nullarg(),
        i.e. the argument is omitted.
        
        This is used for nested calls where a parameter is an
        argument to another call. If the parameter has not been passed an
        argument it is often preferable to similarly omit it in the deeper
        call rather than pass it with its default value. This is particularly
        the case when a parameter gets its mode and type from an argument. If
        it is not passed an argument it has no mode or type and cannot in
        turn be passed as the argument to another call.


    \section{numfields - get the number of fields in a struct}\label{sec:ifnumfields}
        
        {\bf Call:} \\
        numfields(v)
        
        {\bf parameters:} \\
        \begin{tabular}{| l | l | l | l |}
        \hline
        param  & mode & type & \\
        \hline \hline
        v & I V SV S Q & struct & variable  \\ \hline
        \end{tabular}
        
        {\bf Parameter=argument allowed:} \\
        No
        
        {\bf return mode:} \\
        immediate

        {\bf return type:} \\
        int
        
        {\bf Description:} \\
        This function returns the number of fields in the struct given as the
        argument.


    \section{numinargs - number of input arguments}\label{sec:ifnuminargs}
        
        
        {\bf Call:} \\
        numinargs()
        
        {\bf parameters:} \\
        none
        
        {\bf return mode:} \\
        immediate

        {\bf return type:} \\
        int
        
        {\bf Description:} \\
        Returns the number of input arguments to the current module, procedure
        or function.

        See also {\bf numwords()} \autorefph{sec:ifnumwords} which returns the
        number of words in a variable, experssion or type.


    \section{numoutargs - number of output arguments}\label{sec:ifnumoutargs}
        
        
        {\bf Call:} \\
        numoutargs()
        
        {\bf parameters:} \\
        none
        
        {\bf return mode:} \\
        immediate

        {\bf return type:} \\
        int
        
        {\bf Description:} \\
        Returns the number of output arguments to the current module procedure
        or function (returns 0 for a function).


    \section{numwords - get the number of words in a variable, expression or type}\label{sec:ifnumwords}
        
        {\bf Call:} \\
        numwords(v)
        
        {\bf parameters:} \\
        \begin{tabular}{| l | l | l | l |}
        \hline
        param  & mode & type & \\
        \hline \hline
        v & any & any & variable  \\ \hline
        \end{tabular}
        
        {\bf Parameter=argument allowed:} \\
        No
        
        {\bf return mode:} \\
        immediate

        {\bf return type:} \\
        int
        
        {\bf Description:} \\
        This function returns the number of words in the variable or expression
        given as the argument. The argument may be any mode except clock,
        cmemory or rmemory. If the argument has fields or subscripts the number
        of words in the reference is retuned. If the variable is an immediate
        type "type" the number of words in the type is returned.
        
        The number of words is the number of fields or array members in the
        type, including any nested types. In the following example -
        \begin{lstlisting}[gobble=8, language=threepl]
           struct(s,
               f1, "log",
               f2, "[2][3]int:8",
               f3, "bits:4"
           );
           static(ss, s);
           println(numwords(s));
           println(numwords(ss));
           println(numwords(ss.f2));
        \end{lstlisting}
        would print
        \begin{verbatim}
        8
        8
        6
        \end{verbatim}

        See also {\bf numfields()} \autorefph{sec:ifnumfields} which returns the
        number of fields in a struct.
    


    \section{offset - point offset of a fixed point value}\label{sec:ifoffset}
        
        {\bf Call:} \\
        offset(arg)
        
        {\bf parameters:} \\
        \begin{tabular}{| l | l | l | l |}
        \hline
        param  & mode & type & \\
        \hline \hline
        arg  & V SV S Q IN & fixed ufixed & fixed point value  \\ \hline
        \end{tabular}
        
        {\bf Parameter=argument allowed:} \\
        No
        
        {\bf return mode:} \\
        immediate

        {\bf return type:} \\
        uint
        
        {\bf Description:} \\
        This function returns the binary point offset of a fixed point value,
        i.e. the number of bits to the right of the fixed binary point.
 

    \section{ord - get the ordinal of an enum or a character}\label{sec:iford}
        
        {\bf Call:} \\
        ord(v)
        
        {\bf parameters:} \\
        \begin{tabular}{| l | l | l | l |}
        \hline
        param  & mode & type & \\
        \hline \hline
        v & I V SV S Q & enum & enumerated type value or string \\ \hline
        \end{tabular}
        
        {\bf Parameter=argument allowed:} \\
        No
        
        {\bf return mode:} \\
        immediate or value

        {\bf return type:} \\
        int or uint
        
        {\bf Description:} \\
        Returns the ordinal of an enumerated type value or a character.
        The ordinal is a positive integer.
        
        If the argument is immediate mode enum type the returned result is
        immediate mode and the type is uint.
        
        If the argument is a target mode enum type (value, selectvalue, static or
        queue) the returned result is value mode type uint. For an argument which is
        a queue the result will only be available when the queue is available.
        
        If the argument is immediate mode string type it must be of length 1
        (a single character). The ordinal value of that character is returned
        (the ASCII code).


    \section{period - get the period of a clock variable}\label{sec:ifperiod}
        
        {\bf Call:} \\
        period(v)
        
        {\bf parameters:} \\
        \begin{tabular}{| l | l | l | l |}
        \hline
        param  & mode & type & \\
        \hline \hline
        v & C & - & clock variable \\ \hline
        \end{tabular}
        
        {\bf Parameter=argument allowed:} \\
        No
        
        {\bf return mode:} \\
        immediate

        {\bf return type:} \\
        float
        
        {\bf Description:} \\
        Returns the period of a clock variable in ns. If the clock period
        is variable this is the minimum clock period.
        
        See also inbuilt functions - \\
        {\bf frequency()} \autorefph{sec:iffrequency} \\
        {\bf dutycycle()} \autorefph{sec:ifdutycycle}





    \section{portdomain - get the port clock domain for a memory variable}\label{sec:ifportdomain}
        
        {\bf Call:} \\
        readdomain(ident, port)
        
        {\bf parameters:} \\
        \begin{tabular}{| l | l | l | l |}
        \hline
        param  & mode & type & \\
        \hline \hline
        ident & - &  -  & variable identifier \\ \hline
        port  & - &  -  & port number \\ \hline
        \end{tabular}
        
        {\bf Parameter=argument allowed:} \\
        No
        
        {\bf return mode:} \\
        clock

        {\bf return type:} \\
        none
        
        {\bf Description:} \\
        Returns the clock domain for a port of a cmemory or rmemory mode
        variable. If the memory port has not had a read or write procedure
        called thus far then null will be returned.





    \section{pow - raise to a power}\label{sec:ifpow}
           
        {\bf Call:} \\
        pow(f1, f2)
        
        {\bf parameters:} \\
        \begin{tabular}{| l | l | l | l |}
        \hline
        param  & mode & type & \\
        \hline \hline
        f1 & I & float, uint, int & \\ \hline
        f2 & I & float, uint, int & \\ \hline
        \end{tabular}
        
        {\bf Parameter=argument allowed:} \\
        No
        
        {\bf return mode:} \\
        immediate

        {\bf return type:} \\
        float or int
        
        {\bf Description:} \\
        Returns the first argument raised to the power of the second argument.
        Both arguments must be immediate. If an argument is "int" or "uint" type
        it is converted to "float" type. If neither arguments were type "float"
        and the second argument is not negative the result is returned as type
        "int".
 

    \section{prev - get the previous value of an enum}\label{sec:ifprev}
        
        {\bf Call:} \\
        prev(v)
        
        {\bf parameters:} \\
        \begin{tabular}{| l | l | l | l |}
        \hline
        param  & mode & type & \\
        \hline \hline
        v & I V SV S Q & enum & enumerated type value \\ \hline
        \end{tabular}
        
        {\bf Parameter=argument allowed:} \\
        No
        
        {\bf return mode:} \\
        immediate or value

        {\bf return type:} \\
        enum
        {\bf Description:} \\
        Returns the previous value of an enumerated type. The previous value
        is the predecessor in the order the values had as arguments to
        the procedure which created the enum type.
        
        If the argument is immediate mode the returned result is immediate
        mode and the type is the same as the argument. If there is no previous
        value a fatal error occurs.
        
        If the argument is a target mode (value, selectvalue, static or queue)
        the returned result is value mode and the type is the same as the
        argument. For an argument which is a queue the result will only be
        available when the queue is available. If there is no previous
        value a value with an ordinal of 0 is returned.


    \section{prevclock - get the previous clock variable}\label{sec:ifprevclock}
        
        {\bf Call:} \\
        prevclock()
        
        {\bf parameters:} \\
        none
        
        {\bf return mode:} \\
        clock

        {\bf return type:} \\
        none
        
        {\bf Description:} \\
        Gets the previous clock variable, i.e. the one in use prior to
        the most recent {\bf useclock()} call. Null will be returned if there
        was no previous clock domain or the previous clock domain was explicitly
        null.
 

    \section{prielse - the else case value for a priority variable}\label{sec:ifprielse}
        
        {\bf Call:} \\
        prielse(p)
        
        {\bf parameters:} \\
        \begin{tabular}{| l | l | l | l |}
        \hline
        param  & mode & type & \\
        \hline \hline
        p & PR & []log & priority variable identifier \\ \hline
        \end{tabular}
        
        {\bf Parameter=argument allowed:} \\
        No
        
        {\bf return mode:} \\
        value

        {\bf return type:} \\
        log
        
        {\bf Description:} \\
        Return the else case of a priority variable. The returned
        VALUE variable is true if all the priority variable inputs
        are low.



    \section{primtypestring - type of a primitive variable as a string}\label{sec:ifprimtypestring}
        
        {\bf Call:} \\
        primtypestring(ident)
        
        {\bf parameters:} \\
        \begin{tabular}{| l | l | l | l |}
        \hline
        param  & mode & type & \\
        \hline \hline
        ident & - &  -  & variable identifier \\ \hline
        \end{tabular}
        
        {\bf Parameter=argument allowed:} \\
        No
        
        {\bf return mode:} \\
        immediate

        {\bf return type:} \\
        str
        
        {\bf Description:} \\
        Return the type of a primitive variable or expression as a string
        description. If the variable is a target type the colon and width are
        omitted from the string. If the variable is not a primitive type, e.g. it
        is an array or struct, then "none" is returned. If the variable is
        an enum type then the string for that type will be returned.
        
        For modes {\bf cmemory} or {\bf rmemory} "none" is returned. For
        mode {\bf clock} "log" is returned.
        
        Note that if the argument is the identifier for a type variable or
        a variable created by calling struct(), or is the return value from
        inbuilt funtion type(), then
        {\bf primtypestring()} will simply return "type" .


        See also - \\
        {\bf type()} \autorefph{sec:iftype} \\
        {\bf arraytype()} \autorefph{sec:ifarraytype} \\
        {\bf imtotargtype()} \autorefph{sec:ifimtotargtype} \\
        {\bf targtoimtype()} \autorefph{sec:iftargtoimtype} \\
        {\bf typestring()} \autorefph{sec:iftypestring} \\
        {\bf typeval()} \autorefph{sec:iftypeval} \\
        {\bf typevalstring()} \autorefph{sec:iftypevalstr} \\
        {\bf typeisequal()} \autorefph{sec:iftypeiseq} \\
 

    \section{queuedepth - get the depth of a queue}\label{sec:ifqueuedepth}
        
        {\bf Call:} \\
        queuedepth(v)
        
        {\bf parameters:} \\
        \begin{tabular}{| l | l | l | l |}
        \hline
        param  & mode & type & \\
        \hline \hline
        v & Q & - & queue variable \\ \hline
        \end{tabular}
        
        {\bf Parameter=argument allowed:} \\
        No
        
        {\bf return mode:} \\
        immediate

        {\bf return type:} \\
        int
        
        {\bf Description:} \\
        Returns the depth of a queue variable buffer. The depth is set by setting
        attributes {\bf depth} or {\bf mindepth} for the queue variable.
        
        The number of entries in a synchronous queue during target execution can
        be determined by the value returned from inbuilt function {\bf
        queuewords()} (\autorefph{sec:ifqueuewords}). For asynchronous queues the
        library function {\bf asyncqueuewords()}
        (\autorefph{sec:lfasyncqueuewords}) will give a guaranteed minimum
        instead of the actual count. Either of these this can only be called
        from within a module which reads from the queue.
        
        The number of free spaces in a synchronous queue can be determined
        by the value returned from inbuilt function {\bf queuespaces()}
        (\autorefph{sec:ifqueuespaces}). For asynchronous queues the
        library function {\bf asyncqueuespaces()}
        (\autorefph{sec:lfasyncqueuespaces}) will give a guaranteed minimum
        instead of the actual count. Either of these this can only be called
        from within a module which writes to the queue.
 

    \section{queueisempty - determine if a queue is empty}\label{sec:ifqueueisempty}
        
        {\bf Call:} \\
        queueisempty(v)
        
        {\bf parameters:} \\
        \begin{tabular}{| l | l | l | l |}
        \hline
        param  & mode & type & \\
        \hline \hline
        v & Q & - & synchronous queue variable \\ \hline
        \end{tabular}
        
        {\bf Parameter=argument allowed:} \\
        No
        
        {\bf return mode:} \\
        immediate

        {\bf return type:} \\
        log
        
        {\bf Description:} \\
        Returns true if the queue is empty. This function cannot be used with an asynchronous
        queue. The advantage of this function over simply the inverse of the
        read availability (operator {\bf \verb'<'}) is that it can be used
        in a module which does not read from the queue, for example a module
        which writes to a queue can tell if all entries have been read.
        
        The number of entries in a synchronous queue during target execution can
        be determined by the value returned from inbuilt function {\bf
        queuewords()} (\autorefph{sec:ifqueuewords}). For asynchronous queues the
        library function {\bf asyncqueuewords()}
        (\autorefph{sec:lfasyncqueuewords}) will give a guaranteed minimum
        instead of the actual count. Either of these this can only be called
        from within a module which reads from the queue.
        
        The number of free spaces in a synchronous queue can be determined
        by the value returned from inbuilt function {\bf queuespaces()}
        (\autorefph{sec:ifqueuespaces}). For asynchronous queues the
        library function {\bf asyncqueuespaces()}
        (\autorefph{sec:lfasyncqueuespaces}) will give a guaranteed minimum
        instead of the actual count. Either of these this can only be called
        from within a module which writes to the queue.
 

    \section{queuemaxdepth - get the maximum allowed depth of a queue}\label{sec:ifqueuemaxdepth}
        
        {\bf Call:} \\
        queuemaxdepth()
        
        {\bf parameters:} \\
        none
        
        {\bf return mode:} \\
        immediate

        {\bf return type:} \\
        int
        
        {\bf Description:} \\
        Returns the maximum allowed depth of a queue variable buffer.
 

    \section{queuemindepth - get the minimum allowed depth of an asynchronous queue}\label{sec:ifqueuemindepth}
        
        {\bf Call:} \\
        queuemindepth()
        
        {\bf parameters:} \\
        none
        
        {\bf return mode:} \\
        immediate

        {\bf return type:} \\
        int
        
        {\bf Description:} \\
        Returns the minimum allowed depth of an asynchronous queue variable buffer.
        This is also the minimum allowed depth of a synchronous queue variable
        buffer above depth 2 (the default depth).
  

    \section{queuereadstatus - get the read status of an asynchronous
                                                queue}\label{sec:ifqueuereadstatus}
        
        {\bf Call:} \\
        queuereadstatus(v)
        
        {\bf parameters:} \\
        \begin{tabular}{| l | l | l | l |}
        \hline
        param  & mode & type & \\
        \hline \hline
        v & Q & - & queue variable \\ \hline
        \end{tabular}
        
        {\bf Parameter=argument allowed:} \\
        No
        
        {\bf return mode:} \\
        value

        {\bf return type:} \\
        struct
        
        {\bf Description:} \\
        Returns the read status of an asynchronous queue, i.e. a queue which
        is assigned in a different clock domain to that in which it is
        read. The status is a struct with the following fields of
        type "log" -

        \begin{tabular}{| l | l |}
        \hline
        etoq & empty to 1/4 full        \\ \hline
        ttoh & 2 words to 1/2 full      \\ \hline
        qtot & 1/4 full + 2 to 3/4 full \\ \hline
        htof & 1/2 full + 2 to full     \\ \hline 
        ttof & 3/4 full + 2 to full     \\ \hline  
        \end{tabular}

        The five fields are mutually exclusive and one of them is always true.

        {\bf Note that queuereadstatus() is similar to a queue read in that it is
        synchronised to the queue output clock. It must be called in the same
        clock domain as reads of the queue and if the call precedes any reads in
        the source code it will set the read clock domain.}

        The read status has no latency with respect to the read clock. The empty
        condition can be monitored with the read availability operator.
        
        There is a complementary inbuilt function {\bf queuewritestatus()}
        (\autorefph{sec:ifqueuewritestatus}) to get the status in the write
        clock domain.

        Where a queue is synchronous, i.e. assignments to it occur in the same
        clock domain in which it is read, the {\bf queuereadstatus()} function
        cannot be used (this is checked). Instead the {\bf queuewords()}
        (\autorefph{sec:ifqueuewords}) function can be used to get the number of
        buffer entries and {\bf queuespaces()} (\autorefph{sec:ifqueuespaces})
        returns the number of available spaces.
        
        There is a library function {\bf asyncqueuewords()}
        (\autorefph{sec:lfasyncqueuewords}) which can be used to compute
        a guaranteed minimum word count from the read status.
        Similarly there is a library function {\bf asyncqueuespaces()}
        (\autorefph{sec:lfasyncqueuespaces}) which can be used to compute
        a guaranteed minimum space count from the write status.

        For any queue variable {\bf queuedepth()} (\autorefph{sec:ifqueuedepth}) gives
        the depth of the queue buffer.


    \section{queuespaces - get the number of free spaces in a queue}\label{sec:ifqueuespaces}
        
        {\bf Call:} \\
        queuespaces(v)
        
        {\bf parameters:} \\
        \begin{tabular}{| l | l | l | l |}
        \hline
        param  & mode & type & \\
        \hline \hline
        v & Q & - & queue variable \\ \hline
        \end{tabular}
        
        {\bf Parameter=argument allowed:} \\
        No
        
        {\bf return mode:} \\
        value

        {\bf return type:} \\
        uint
        
        {\bf Description:} \\
        Returns the number of free entries in the buffer of a synchronous queue,
        i.e. a queue which is assigned in the same clock domain to that in which
        it was created.

        A call to {\bf queuespaces()} must be in the same clock domain as all
        reads and writes of the queue. If {\bf queuespaces()} is called prior to
        any queue reads or writes it will set the queue clock domain to the
        current domain.

        For an asynchronous queue, i.e. a queue which is assigned in a different
        clock domain to that in which it is read, the occupancy status can
        only be approximately determined by using inbuilt functions {\bf
        queuereadstatus()} (\autorefph{sec:ifqueuereadstatus}) and {\bf
        queuewritestatus()} (\autorefph{sec:ifqueuewritestatus}).
        There is a library function {\bf asyncqueuespaces()}
        (\autorefph{sec:lfasyncqueuespaces}) which can be used to compute
        a guaranteed minimum space count from the write status.
        Similarly there is a library function {\bf asyncqueuewords()}
        (\autorefph{sec:lfasyncqueuewords}) which can be used to compute
        a guaranteed minimum word count from the read status.

        To get the number of occupied entries use {\bf queuewords()} (section
        \autorefph{sec:ifqueuewords}).
        
        See \autorefph{sec:attributes} (queue mode
        attributes) for details on the depth of queue buffers. To get the queue
        buffer depth use {\bf queuedepth()} (\autorefph{sec:ifqueuedepth}).
        
        For a synchronous queue the inbuilt function {\bf queueisempty()}
        (\autorefph{sec:ifqueueisempty}) can be called. Unlike the read
        availability operator this can be used in a module which does not
        read from the queue.
      
 

    \section{queuewords - get the number of entries in a queue}\label{sec:ifqueuewords}
        
        {\bf Call:} \\
        queuewords(v)
        
        {\bf parameters:} \\
        \begin{tabular}{| l | l | l | l |}
        \hline
        param  & mode & type & \\
        \hline \hline
        v & Q & - & queue variable \\ \hline
        \end{tabular}
        
        {\bf Parameter=argument allowed:} \\
        No
        
        {\bf return mode:} \\
        value

        {\bf return type:} \\
        uint
        
        {\bf Description:} \\
        Returns the number of occupied entries in the buffer of a synchronous
        queue, i.e. a queue which is assigned in the same clock domain to that in
        which it was created.

        A call to {\bf queuewords()} must be in the same clock domain as all
        reads and writes of the queue. If {\bf queuewords()} is called prior to
        any queue reads or writes it will set the queue clock domain to the
        current domain.

        For an asynchronous queue, i.e. a queue which is assigned in a different
        clock domain to that in which it is read, the occupancy status can only
        be approximately determined by using inbuilt functions {\bf
        queuereadstatus()} (\autorefph{sec:ifqueuereadstatus}) and {\bf
        queuewritestatus()} (\autorefph{sec:ifqueuewritestatus}).
        There is a library function {\bf asyncqueuewords()}
        (\autorefph{sec:lfasyncqueuewords}) which can be used to compute
        a guaranteed minimum word count from the read status.
        Similarly there is a library function {\bf asyncqueuespaces()}
        (\autorefph{sec:lfasyncqueuespaces}) which can be used to compute
        a guaranteed minimum space count from the write status.

        To get the number of free entries use {\bf queuespaces()} (section
        \autorefph{sec:ifqueuespaces}).
        
        See \autorefph{sec:attributes} (queue mode attributes) for details on the
        depth of queue buffers. To get the queue buffer depth use {\bf
        queuedepth()} (\autorefph{sec:ifqueuedepth}).
   

    \section{queuewritestatus - get the write status of an asynchronous
                                                queue}\label{sec:ifqueuewritestatus}
        
        {\bf Call:} \\
        queuewritestatus(v)
        
        {\bf parameters:} \\
        \begin{tabular}{| l | l | l | l |}
        \hline
        param  & mode & type & \\
        \hline \hline
        v & Q & - & queue variable \\ \hline
        \end{tabular}
        
        {\bf Parameter=argument allowed:} \\
        No
        
        {\bf return mode:} \\
        value

        {\bf return type:} \\
        struct
        
        {\bf Description:} \\
        Returns the read status of an asynchronous queue, i.e. a queue which
        is assigned in a different clock domain to that in which it is
        read. The status is a struct with the following fields of
        type "log" -

        \begin{tabular}{| l | l |}
        \hline
        etoq & empty to 1/4 full        \\ \hline
        ttoh & 2 words to 1/2 full      \\ \hline
        qtot & 1/4 full + 2 to 3/4 full \\ \hline
        htof & 1/2 full + 2 to full     \\ \hline 
        ttof & 3/4 full + 2 to full     \\ \hline  
        \end{tabular}

        {\bf Parameter=argument allowed:} \\
        No
        
        The five fields are mutually exclusive and one of them is always true.
        
        {\bf Note that queuewritestatus() is similar to an assignment in that it
        is synchronised to the queue input clock. It must be called in the same
        clock domain as assignments to the queue and if the call precedes any
        assignments in the source code it will set the assignment clock domain.}

        The write status has no latency with respect to the write clock. The
        full condition can be monitored using the write availability operator.
        
        There is a complementary inbuilt function {\bf queuereadstatus()}
        (\autorefph{sec:ifqueuereadstatus}) to get the status in the read
        clock domain.

        Where a queue is synchronous, i.e. assignments to it occur in the same
        clock domain in which it was created the {\bf queuewritestatus()}
        function cannot be used (this is checked). Instead the {\bf
        queuewords()} (\autorefph{sec:ifqueuewords}) function can be used to get the number
        of buffer entries and {\bf queuespaces()} (\autorefph{sec:ifqueuespaces})
        returns the number of available spaces.
        
        There is a library function {\bf asyncqueuespaces()}
        (\autorefph{sec:lfasyncqueuespaces}) which can be used to compute
        a guaranteed minimum space count from the write status.
        Similarly there is a library function {\bf asyncqueuewords()}
        (\autorefph{sec:lfasyncqueuewords}) which can be used to compute
        a guaranteed minimum word count from the read status.

        For any queue variable {\bf queuedepth()} (\autorefph{sec:ifqueuedepth}) gives
        the depth of the queue buffer.



    \section{random - a random number}\label{sec:ifrandom}
            
        {\bf Call:} \\
        random()
        
        {\bf parameters:} \\
        none.
        
        {\bf return mode:} \\
        immediate

        {\bf return type:} \\
        float
        
        {\bf Description:} \\
        Returns an immediate random number. The random number {\em n} is in
        the range $0.0 \le n < 1.0$.


    \section{readdomain - get the read clock domain for a variable}\label{sec:ifreaddomain}
        
        {\bf Call:} \\
        readdomain(ident)
        
        {\bf parameters:} \\
        \begin{tabular}{| l | l | l | l |}
        \hline
        param  & mode & type & \\
        \hline \hline
        ident & - &  -  & variable identifier \\ \hline
        \end{tabular}
        
        {\bf Parameter=argument allowed:} \\
        No
        
        {\bf return mode:} \\
        clock

        {\bf return type:} \\
        none
        
        {\bf Description:} \\
        Returns the read (output) clock domain for a variable. If the variable
        has not had its value referenced so far null will be returned as the
        variable will not yet have had its output clock determined. The
        variable may be selectvalue, value, static, queue or priority mode.
        It cannot be memory or clock mode.
        
        See also {\bf writedomain()} \autorefph{sec:ifwritedomain} which gets the write
        (input or assignment) clock and {\bf domain()} \autorefph{sec:ifdomain} which
        gets the single clock for selectvalue, value, static and priority mode
        variables.
        
        For a queue the read and write clocks may be the same or may differ.
        For selectvalue, value, static and priority mode variables there is only
        one clock and all three inbuilt functions will return that clock.


    \section{readbyte - read a byte from an input file}\label{sec:ifreadbyte}
        
        {\bf Call:} \\
        readbyte(file)
        
        {\bf parameters:} \\
        \begin{tabular}{| l | l | l | l |}
        \hline
        param  & mode & type & \\
        \hline \hline
        file & I & file & open input file  \\ \hline
        \end{tabular}
        
        {\bf Parameter=argument allowed:} \\
        No
        
        {\bf return mode:} \\
        immediate

        {\bf return type:} \\
        int
        
        {\bf Description:} \\
        This function reads a byte from an input file. The argument is
        an open input file variable. If a read beyond the end-of-file is
        attempted -1 is returned and the inbuilt function {\bf
        hadeof()} will return true.
        
        See the note on character and byte reads and writes in the description
        of inbuilt procedure {\bf file()} \autorefph{sec:ipfile}.
        
        See also - \\
        {\bf scanf()} \autorefph{sec:ipscanf} \\
        {\bf fscanf()} \autorefph{sec:ipfscanf} \\
        {\bf sscanf()} \autorefph{sec:ipsscanf} \\
        {\bf readln()} \autorefph{sec:ifreadln}.


    \section{readln - read a line from an input file}\label{sec:ifreadln}
        
        {\bf Call:} \\
        readln(file)
        
        {\bf parameters:} \\
        \begin{tabular}{| l | l | l | l |}
        \hline
        param  & mode & type & \\
        \hline \hline
        file & I & file & open input file  \\ \hline
        \end{tabular}
        
        {\bf Parameter=argument allowed:} \\
        No
        
        {\bf return mode:} \\
        immediate

        {\bf return type:} \\
        str
        
        {\bf Description:} \\
        This function reads a line of text from an input file. The argument is
        an open input file variable. If a read beyond the end-of-file is
        attempted an empty string is returned and the inbuilt function {\bf
        hadeof()} will return true.       
        
        See the note on character and byte reads and writes in the description
        of inbuilt procedure {\bf file()} \autorefph{sec:ipfile}.
        
        See also - \\
        {\bf scanf()} \autorefph{sec:ipscanf} \\
        {\bf fscanf()} \autorefph{sec:ipfscanf} \\
        {\bf sscanf()} \autorefph{sec:ipsscanf} \\
        {\bf readbyte()} \autorefph{sec:ifreadbyte}.
       

    \section{replace - replace left substring}\label{sec:ifreplace}
        
        {\bf Call:} \\
        replace(s1, s2, s3)
        
        {\bf parameters:} \\
        \begin{tabular}{| l | l | l | l |}
        \hline
        param  & mode & type & \\
        \hline \hline
        s1 & I & str & subject string       \\ \hline
        s2 & I & str & selection string     \\ \hline
        s3 & I & str & replacement string   \\ \hline
        \end{tabular}
        
        {\bf Parameter=argument allowed:} \\
        No
        
        {\bf return mode:} \\
        immediate

        {\bf return type:} \\
        str
        
        {\bf Description:} \\
        Returns the string supplied by the first argument with the leftmost
        substring matched by the second argument replaced by the third argument
        string. If the substring is not found the first argument is returned
        unchanged.
        
        The second argument string is a Java regular expression - 
        see \autorefph{chap:regexp}.


    \section{replaceall - replace all substrings}\label{sec:ifreplaceall}
        
        {\bf Call:} \\
        replaceall(s1, s2, s3)
        
        {\bf parameters:} \\
        \begin{tabular}{| l | l | l | l |}
        \hline
        param  & mode & type & \\
        \hline \hline
        s1 & I & str & subject string       \\ \hline
        s2 & I & str & selection string     \\ \hline
        s3 & I & str & replacement string   \\ \hline
        \end{tabular}
        
        {\bf Parameter=argument allowed:} \\
        No
        
        {\bf return mode:} \\
        immediate

        {\bf return type:} \\
        str
        
        {\bf Description:} \\
        Returns the string supplied by the first argument with all
        substrings matched by the second argument replaced by the third argument
        string. If the substring is not found the first argument is returned
        unchanged.
        
        The second argument string is a Java regular expression - 
        see \autorefph{chap:regexp}.
      
 

    \section{resetoutput - get the reset output of a queue}\label{sec:ifresetoutput}
        
        {\bf Call:} \\
        resetoutput(v)
        
        {\bf parameters:} \\
        \begin{tabular}{| l | l | l | l |}
        \hline
        param  & mode & type & \\
        \hline \hline
        v & Q & - & queue variable \\ \hline
        \end{tabular}
        
        {\bf Parameter=argument allowed:} \\
        No
        
        {\bf return mode:} \\
        value

        {\bf return type:} \\
        log
        
        {\bf Description:} \\
        Returns a value of type log which is true for one clock cycle when the
        queue is reset.

        {\bf Note that resetoutput() is similar to a queue read in that it is
        synchronised to the queue output clock. It must be called in the same
        clock domain as reads of the queue and if the call precedes any reads in
        the source code it will set the read clock domain.}

        For an asynchronous queue the reset output is synchronised to the output
        clock (read clock domain) and follows the reset input to the queue
        which is synchronised to the input clock (write clock domain).
        
        The purpose of this function is to allow queue resets of pipelines with
        multiple clock domains to be chained from the source end of a pipeline.
        The source end of a pipeline is explicitly reset and each clock domain
        section of the pipeline downstream derives its reset signal from the
        asynchronous queue reset output at the head of the section.
 


    \section{round - round float or fixed to nearest integer}\label{sec:ifround}
        
        
        {\bf Call:} \\
        round(f)
        
        {\bf parameters:} \\
        \begin{tabular}{| l | l | l | l |}
        \hline
        param  & mode & type & \\
        \hline \hline
        f & I & float & \\ \hline
          & V & fixed, ufixed & \\ \hline
        \end{tabular}
        
        {\bf Parameter=argument allowed:} \\
        No
        
        {\bf return mode:} \\
        immediate or value

        {\bf return type:} \\
        int or uint
        
        {\bf Description:} \\
        For an immediate argument of type float this rounds the floating point
        argument to the nearest integer. An int type is returned.
        
        For a value argument of type fixed or ufixed this rounds the
        fixed point argument using round-to-even. This method rounds a
        fraction which is exactly $1/2$ to the nearest even integer to avoid
        a statistical bias. Type int is returned for fixed and uint for ufixed.
        
        Note that rounding of an integer target type at a selected bit position
        can be accomplished by first casting it to a fixed or ufixed type with a
        fractional part of appropriate size then applying the round() function.



    \section{rmemoryout - get the value of a registered output memory}\label{sec:ifrmemoryout}
       
        {\bf Call:} \\
        rmemoryout(id, port)
        
        {\bf parameters:} \\
        \begin{tabular}{| l | l | l | l |}
        \hline
        param  & mode & type & \\
        \hline \hline
        id   & M &  -  & memory identifier  \\ \hline
        port & I & int & port               \\ \hline
        \end{tabular}
        
        {\bf Parameter=argument allowed:} \\
        Yes
        
        {\bf return mode:} \\
        value

        {\bf return type:} \\
        memory data type
        
        {\bf Description:} \\
        Get a value which is the output of a registered memory port.
        This value is continuously valid. The type is the same as the
        data type argument provided when the memory was created.
        
        A call to {\bf rmemoryout()} sets the clock domain for the memory port to
        the current clock domain. All calls to {\bf rmemoryout()}
        (\autorefph{sec:ifrmemoryout}), {\bf rmemoryread()} (\autorefph{sec:iprmemoryread}) or {\bf
        rmemorywrite()} (\autorefph{sec:iprmemorywrite}) for the same port of the same
        registered output memory must be in the same clock domain (or a fatal
        error results).
            


    \section{sample - sample an expression}\label{sec:ifsample}

        {\bf Call:} \\
        sample(e)

        {\bf parameters:} \\
        \begin{tabular}{| l | l | l | l |}
        \hline
        param  & mode & type & \\
        \hline \hline
        e & - & any & target expression  \\ \hline
        \end{tabular}

        {\bf Parameter=argument allowed:} \\
        No
        
        {\bf return mode:} \\
        value

        {\bf return type:} \\
        type of argument

        {\bf Description:} \\
        The argument must be a static variable or a value (an expression or a
        value variable) containing static variables or a value variable which
        has been passed as an output argument to the {\bf import()} inbuilt
        procedure. The argument must not contain queue reads and cannot be a
        memory or clock mode variable.

        The argument value will be sampled in the current clock domain and the
        resulting value returned as the value of the function. The argument must
        have a clock domain or a fatal error will occur. It will have one if has
        already been evaluated previously (in a clock domain). If evaluation
        occurs later in the source code then the clock domain of the argument
        must have been explicitly forced prior to calling {\bf sample()}. This
        can be done by using inbuilt procedure {\bf domain()} or {\bf
        readdomain()} on any variables in the argument. The argument cannot be a
        constant.

        The simplest case where this is used is where a static variable
        in one clock domain is assigned to a static variable (or a queue)
        in another clock domain. The sampling avoids problems
        with violations of setup and hold time which might otherwise be
        encountered. If the two clocks have an exact frequency ratio
        with a fixed phase then sampling may not be required as
        long as an appropriate clock-to-clock timing constraint is
        imposed.
        \index{constraint!timing}
        \index{timing constraint}
        If the source variable changes very seldom and an
        occasional transient data error can be tolerated then again
        sampling can be omitted (see inbuilt function {\bf accept()}
        \autorefph{sec:ifaccept}).

        This is a simple device to avoid the use of a full dual-port
        resynchronising queue buffer where the data is known to be sparse or
        where the data is sampled rather than pipelined. If the input clock
        frequency is the higher of the two then some data will not be sampled.
        If the input clock frequency is the lower of the two then data will be
        multiply sampled. If the two clock frequencies are very close both missing
        samples and duplicated samples may occur. Note that
        regardless of missing or duplicated samples the sampling process
        maintains data order, i.e. data are never out of sequence.

        A typical use of this function is when the source static
        variable seldom changes, for example a controlling
        parameter assigned via the CPU bus, but the value is used
        in another clock domain to the interface static variable.
  
        {\bf sample()} should not be used for detecting single cycle pulses in
        another clock domain - in that case see inbuilt module {\bf resync()}
        \autorefph{sec:imresync} or library function {\bf resync()} \autorefph{sec:lfresync} which
        guarantee a single cycle pulse regardless of relative clock frequencies.

        The sampled data may be delayed up to two clock cycles in
        the current clock domain.
        
        If the argument clock domain is the same as the current clock domain
        the argument is simply returned as the value.



    \section{selecta - a combinatorial selector from arrays}\label{sec:ifselecta}

        {\bf Call:} \\
        selecta(sel, data)
        
        {\bf parameters:} \\
        \begin{tabular}{| l | l | l | l |}
        \hline
        param  & mode & type & \\
        \hline \hline
        sel  & V & []log & selector array \\ \hline
        data & V & []any & data array     \\ \hline
        \end{tabular}
        
        {\bf Parameter=argument allowed:} \\
        No
        
        {\bf return mode:} \\
        value

        {\bf return type:} \\
        type of data inputs
        
        {\bf Description:} \\
        An inbuilt selector function using gate logic.
        This has 2 arguments. The 1st argument is an array
        of type {\bf log}. The 2nd argument is an array of any type
        whose 1st or only dimension is the same as that of the 1st
        argument array.
        
        Each member of the selector array is a logical value which selects
        the corresponding data element of the data array.
        
        The function has a defined value when at most one select input
        is available and true and the associated data input is available,
        the value reflecting the selected data input. If more than one
        input is selected the value is undefined. If no inputs are
        selected, the value is zero or false depending on data type(s).
        
        If none of the function inputs contain queues, the function value is
        always available. If any function inputs contain queues the function
        value is available when an input selector and its associated data
        input are available and the selector is true. Note that if some
        inputs contain queues but one input selector and its associated data
        do not, the function will be available when that selector is true.
        Note that the use of this function as a gate with a clear value
        when no selectors are true is only possible if the inputs contain
        no queues since any queues in any inputs will render the function
        value unavailable when no input is selected.
        
            
        See also library function {\bf select()} \autorefph{sec:lfselect}
        where the inputs are select-data pairs.



    \section{signum - sign function}\label{sec:ifsignum}
        
        
        {\bf Call:} \\
        signum(v)
        
        {\bf parameters:} \\
        \begin{tabular}{| l | l | l | l |}
        \hline
        param  & mode & type & \\
        \hline \hline
        v & I & int or float & \\ \hline
        \end{tabular}
        
        {\bf Parameter=argument allowed:} \\
        No
        
        {\bf return mode:} \\
        immediate

        {\bf return type:} \\
        int or float as per argument
        
        {\bf Description:} \\
        Returns 0 or 0.0 if the argument is zero, 1 or 1.0 if the argument is +ve,
        or -1 or -1.0 if the argument is -ve..


    \section{sin - trigonometric sine function}\label{sec:ifsin}
        
        
        {\bf Call:} \\
        sin(f)
        
        {\bf parameters:} \\
        \begin{tabular}{| l | l | l | l |}
        \hline
        param  & mode & type & \\
        \hline \hline
        f & I & float & \\ \hline
        \end{tabular}
        
        {\bf Parameter=argument allowed:} \\
        No
        
        {\bf return mode:} \\
        immediate

        {\bf return type:} \\
        float
        
        {\bf Description:} \\
        Returns the sine of the immediate argument. The argument is the
        angle in radian.


    \section{size - get the size in bits}\label{sec:ifsize}
        
        {\bf Call:} \\
        size(ve)
        
        {\bf parameters:} \\
        \begin{tabular}{| l | l | l | l |}
        \hline
        param  & mode & type & \\
        \hline \hline
        ve & V SV S Q PR & -   & variable or expression \\ \hline
        \end{tabular}
        
        {\bf Parameter=argument allowed:} \\
        No
        
        {\bf return mode:} \\
        immediate

        {\bf return type:} \\
        int
        
        {\bf Description:} \\
        Returns the size in bits.
        
        The argument cannot be cmemory, rmemory or clock mode.
        
        The argument type cannot be type {\bf str}, {\bf \verb'->'},
        {\bf file} or immediate {\bf enum}.
        
        If the argument is a target mode (other than memory or clock mode) the size
        in bits will be returned. If the type is primitive it will return the size
        of the primitive type in bits. If the type is structured it will return the
        total size of the type in bits.
        
        If the argument is immediate mode -
        \blist
        \item   int or uint - the minimum number of bits required to represent the
                primitive type is returned (for a +ve integer no allowance is made
                for a sign bit!) but for arguments 0 or 1, 1 is returned and for
                argument -1, 2 is returned
        \item   log - 1 is returned
        \item   float - 64 is returned
        \item   type - the size of the type value is returned (the type must
                be target type(s))
        \blend
        
        In the case of a structured type the argument may have fields or
        subscripts in which case the size of the field or array member will
        be returned.
        
        The argument may be type "type" in which case the size of the type will
        be returned. If the type is a structured type then fields or subscripts
        may be appended and the size of these fields or array members will be
        returned. If the structured type is immediate rather than target then the
        size returned will be 1 for each element since sizes of immediates
        can only be derived from values, not types.


    \section{span - span a head string selected by characters in a set}\label{sec:ifspan}
        
        {\bf Call:} \\
        span(s1, s2)
        
        {\bf parameters:} \\
        \begin{tabular}{| l | l | l | l |}
        \hline
        param  & mode & type & \\
        \hline \hline
        s1 & I & str & subject string   \\ \hline
        s2 & I & str & search string    \\ \hline
        \end{tabular}
        
        {\bf Parameter=argument allowed:} \\
        No
        
        {\bf return mode:} \\
        immediate

        {\bf return type:} \\
        int
        
        {\bf Description:} \\
        Returns the length of the head string from the first argument whose characters
        are contained in the set of characters given by the second argument
        string.
       

    \section{split - split a string into substrings}\label{sec:ifsplit}
        
        {\bf Call:} \\
        split(s1, s2)
        
        {\bf parameters:} \\
        \begin{tabular}{| l | l | l | l |}
        \hline
        param  & mode & type & \\
        \hline \hline
        s1 & I & str & subject string                           \\ \hline
        s2 & I & str & split point regular expression string    \\ \hline
        \end{tabular}
        
        {\bf Parameter=argument allowed:} \\
        No
        
        {\bf return mode:} \\
        immediate

        {\bf return type:} \\
        \verb'[]str'
        
        {\bf Description:} \\
        This function splits a string into substrings, returning them as an array
        of strings. The split point is indicated by a regular expression, supplied
        as the second argument. The substrings within the first argument are
        separated by sections which match the regular expression. If the regular
        expression does not match any section of the subject string then the returned
        array has a dimension of one, the single entry containing the whole subject
        string.     
        
        The second argument string is a Java regular expression - 
        see \autorefph{chap:regexp}.


    \section{sqrt - an immediate square root}\label{sec:ifsqrt}
            
        {\bf Call:} \\
        sqrt()
        
        {\bf parameters:} \\
        \begin{tabular}{| l | l | l | l |}
        \hline
        param  & mode & type & \\
        \hline \hline
        f & I & float & \\ \hline
        \end{tabular}
        
        {\bf Parameter=argument allowed:} \\
        No
        
        {\bf return mode:} \\
        immediate

        {\bf return type:} \\
        float
        
        {\bf Description:} \\
        Returns the square root of the immediate argument.
 

    \section{stacktrace - get the current stack trace}\label{sec:ifstacktrace}
        
        {\bf Call:} \\
        stacktrace()
        
        {\bf parameters:} \\
        none
        
        {\bf return mode:} \\
        immediate

        {\bf return type:} \\
        str
        
        {\bf Description:} \\
        Returns a multi-line string listing the source locations of
        all module, procedure and function calls in the stack.
        
        See also inbuilt procedure {\bf trace()} \autorefph{sec:iptrace}
        which prints the location and stack trace to standard out.


    \section{startswith - a string is the head of another}\label{sec:ifstartswith}
        
        {\bf Call:} \\
        startswith(s1, s2)
        
        {\bf parameters:} \\
        \begin{tabular}{| l | l | l | l |}
        \hline
        param  & mode & type & \\
        \hline \hline
        s1 & I & str & subject string   \\ \hline
        s2 & I & str & test string      \\ \hline
        \end{tabular}
        
        {\bf Parameter=argument allowed:} \\
        No
        
        {\bf return mode:} \\
        immediate

        {\bf return type:} \\
        log
        
        {\bf Description:} \\
        Returns true if the 2nd argument string is the head
        of the 1st argument string.


    \section{strtoarray - convert a string to an array of integers}\label{sec:ifstrtoarray}
        
        {\bf Call:} \\
        strtoarray(s)
        
        {\bf parameters:} \\
        \begin{tabular}{| l | l | l | l |}
        \hline
        param  & mode & type & \\
        \hline \hline
        s & I & str & subject string   \\ \hline
        \end{tabular}
        
        {\bf Parameter=argument allowed:} \\
        No
        
        {\bf return mode:} \\
        immediate

        {\bf return type:} \\
        \verb'[]uint'
        
        {\bf Description:} \\
        Returns an array of integers whose values are the ASCII codes
        representing the argument string.
        
        Note that this function is usually used as an initialiser for an array
        variable and this must be type "[]uint" rather than "[]int", e.g.
        \begin{lstlisting}[gobble=8, language=threepl]
           var(a, "[]uint", strtoarray(s));
        \end{lstlisting}
        This strict typing in this context may be relaxed in the future.

        See also - \\
        \autorefph{sec:ifarraytostr} {\bf arraytostr()} converting
        an integer array to a string.


    \section{strlen - find the length of a string}\label{sec:ifstrlen}
        
        {\bf Call:} \\
        strlen(s)
        
        {\bf parameters:} \\
        \begin{tabular}{| l | l | l | l |}
        \hline
        param  & mode & type & \\
        \hline \hline
        s & I & str & subject string   \\ \hline
        \end{tabular}
        
        {\bf Parameter=argument allowed:} \\
        No
        
        {\bf return mode:} \\
        immediate

        {\bf return type:} \\
        int
        
        {\bf Description:} \\
        Returns the length the argument string.


    \section{strpos - find leftmost substring position}\label{sec:ifstrpos}
        
        {\bf Call:} \\
        strpos(s1, s2)
        
        {\bf parameters:} \\
        \begin{tabular}{| l | l | l | l |}
        \hline
        param  & mode & type & \\
        \hline \hline
        s1 & I & str & subject string   \\ \hline
        s2 & I & str & search string    \\ \hline
        \end{tabular}
        
        {\bf Parameter=argument allowed:} \\
        No
        
        {\bf return mode:} \\
        immediate

        {\bf return type:} \\
        int
        
        {\bf Description:} \\
        Returns the leftmost position of an occurrence of the second argument
        string in the first argument string. If not found, -1 is returned.


    \section{strrpos - find rightmost substring position}\label{sec:ifstrrpos}
        
        {\bf Call:} \\
        strrpos(s1, s2)
        
        {\bf parameters:} \\
        \begin{tabular}{| l | l | l | l |}
        \hline
        param  & mode & type & \\
        \hline \hline
        s1 & I & str & subject string   \\ \hline
        s2 & I & str & search string    \\ \hline
        \end{tabular}
        
        {\bf Parameter=argument allowed:} \\
        No
        
        {\bf return mode:} \\
        immediate

        {\bf return type:} \\
        int
        
        {\bf Description:} \\
        Returns the rightmost position of an occurrence of the second argument
        string in the first argument string. If not found, -1 is returned.


    \section{strtofloat - convert a string to a float}\label{sec:ifstof}
        
        {\bf Call:} \\
        strtofloat(s)
        
        {\bf parameters:} \\
        \begin{tabular}{| l | l | l | l |}
        \hline
        param  & mode & type & \\
        \hline \hline
        s & I & str & subject string   \\ \hline
        \end{tabular}
        
        {\bf Parameter=argument allowed:} \\
        No
        
        {\bf return mode:} \\
        immediate

        {\bf return type:} \\
        float
        
        {\bf Description:} \\
        Returns the floating point value represented by the string.


    \section{strtoint - convert a string to an integer}\label{sec:ifstoi}
        
        {\bf Call:} \\
        strtoint(s)
        
        {\bf parameters:} \\
        \begin{tabular}{| l | l | l | l |}
        \hline
        param  & mode & type & \\
        \hline \hline
        s & I & str & subject string   \\ \hline
        \end{tabular}
        
        {\bf Parameter=argument allowed:} \\
        No
        
        {\bf return mode:} \\
        immediate

        {\bf return type:} \\
        int
        
        {\bf Description:} \\
        Returns the integer value represented by the string. The string
        may begin with 0x or 0X for hexadecimal, 0b or 0B for binary
        or 0 for octal.




    \section{strtotype - convert a string to an type}\label{sec:ifstot}
        
        {\bf Call:} \\
        strtotype(s)
        
        {\bf parameters:} \\
        \begin{tabular}{| l | l | l | l |}
        \hline
        param  & mode & type & \\
        \hline \hline
        s & I & str & subject string   \\ \hline
        \end{tabular}
        
        {\bf Parameter=argument allowed:} \\
        No
        
        {\bf return mode:} \\
        immediate

        {\bf return type:} \\
        type
        
        {\bf Description:} \\
        Returns the type value represented by the string.
        
        A string describing a type can be assigned to a type variable or passed as
        an argument to a parameter which is a type variable. The type string
        may contain the identifier of another type variable, but if the string
        is used in a scope in which the embedded identifier is not
        accessible then the conversion of string-to-type will fail. This
        function converts a string to a type, evaluating an embedded
        string variable in the process in the current scope after which the type
        can be used remote from that embedded variable.




    \section{substitute - substitute characters}\label{sec:ifsubstitute}
        
        {\bf Call:} \\
        substitute(s1, s2, s3)
        
        {\bf parameters:} \\
        \begin{tabular}{| l | l | l | l |}
        \hline
        param  & mode & type & \\
        \hline \hline
        s1 & I & str & subject string           \\ \hline
        s2 & I & str & selection character set  \\ \hline
        s3 & I & str & substitute character set \\ \hline
        \end{tabular}
        
        {\bf Parameter=argument allowed:} \\
        No
        
        {\bf return mode:} \\
        immediate

        {\bf return type:} \\
        str
        
        {\bf Description:} \\
        Returns the string supplied by the first argument with selected
        characters replaced. The second argument is a string whose
        characters constitute the set of selection characters. The third
        argument is a string  whose characters constitute the matching
        set of  replacement characters. If the 3rd argument string is
        shorter than the 2nd then those trailing characters in the 1st
        set will be eliminated. If the 3rd argument string is longer than
        the 2nd then those trailing characters in the 2nd set will be
        ignored. If a character is repeated in the 2nd argument string
        the left-most substitution will occur.                


    \section{substring - get a substring}\label{sec:ifsubstring}
        
        {\bf Call:} \\
        substring(s, i1, i2)
        
        {\bf parameters:} \\
        \begin{tabular}{| l | l | l | l |}
        \hline
        param  & mode & type & \\
        \hline \hline
        s  & I & str & subject string               \\ \hline
        i1 & I & int & starting position            \\ \hline
        i2 & I & int & optional substring length    \\ \hline
        \end{tabular}
        
        {\bf Parameter=argument allowed:} \\
        No
        
        {\bf return mode:} \\
        immediate

        {\bf return type:} \\
        str
        
        {\bf Description:} \\
        Returns a substring of the first argument. The second argument gives
        the starting position of the substring and the optional third argument
        gives its length. If the third argument is missing the substring
        extends to the end of the subject string. If the starting position is
        not within the subject string or the length extends beyond the subjects
        string a fatal error occurs.


    \section{tail - get the tail of a string}\label{sec:iftail}
        
        {\bf Call:} \\
        tail(s1, i2)
        
        {\bf parameters:} \\
        \begin{tabular}{| l | l | l | l |}
        \hline
        param  & mode & type & \\
        \hline \hline
        s1 & I & str & subject string   \\ \hline
        i2 & I & int & length of tail   \\ \hline
        \end{tabular}
        
        {\bf Parameter=argument allowed:} \\
        No
        
        {\bf return mode:} \\
        immediate

        {\bf return type:} \\
        str
        
        {\bf Description:} \\
        Returns the tail of the string first argument, the length of the
        tail string being given by the integer second argument. If the
        length requested is greater than the string length a fatal error
        occurs.


    \section{tan - trigonometric tangent function}\label{sec:iftan}
        
        
        {\bf Call:} \\
        tan(f)
        
        {\bf parameters:} \\
        \begin{tabular}{| l | l | l | l |}
        \hline
        param  & mode & type & \\
        \hline \hline
        f & I & float & \\ \hline
        \end{tabular}
        
        {\bf Parameter=argument allowed:} \\
        No
        
        {\bf return mode:} \\
        immediate

        {\bf return type:} \\
        float
        
        {\bf Description:} \\
        Returns the tangent of the immediate argument. The argument is the
        angle in radian.


    \section{targconsttoimmed - convert a value mode constant to an immediate}\label{sec:iftargconsttoimmed}
        
        {\bf Call:} \\
        targconsttoimmed(v)
        
        {\bf parameters:} \\
        \begin{tabular}{| l | l | l | l |}
        \hline
        param  & mode & type & \\
        \hline \hline
        v & V & prim & value mode variable \\ \hline
        \end{tabular}
        
        {\bf Parameter=argument allowed:} \\
        No
        
        {\bf return mode:} \\
        immediate

        {\bf return type:} \\
        bits, uint, int or log
        
        {\bf Description:} \\
        Return the immediate value of a value variable which is a target
        constant.



    \section{targtoimtype - convert a target type to an immediate type}\label{sec:iftargtoimtype}
        
        {\bf Call:} \\
        targtoimtype(expr)
        
        {\bf parameters:} \\
        \begin{tabular}{| l | l | l | l |}
        \hline
        param  & mode & type & \\
        \hline \hline
        expr & any & any & value \\ \hline
        \end{tabular}
        
        {\bf Parameter=argument allowed:} \\
        No
        
        {\bf return mode:} \\
        immediate

        {\bf return type:} \\
        str
        
        {\bf Description:} \\
        Return the type of a target variable as an immediate type. The type is a
        string in the same format as is used for type declarations. The type
        will be converted to immediate type by removing width fields and converting
        "bits" to "uint".
        
        The argument may be -
        \blist
        \item   a variable of type "type" - the type it represents will
                be converted to an immediate equivalent and returned
                as a string
        \item   a string - the string will be converted to an immediate
                equivalent type string and returned
        \item   a variable - the type of the variable will be converted
                to an immediate equivalent and returned as a string
        \blend
        
        Where the type concerned is already an immediate type it will
        be returned unchanged.
        
        See also - \\
        {\bf type()} \autorefph{sec:iftype} \\
        {\bf arraytype()} \autorefph{sec:ifarraytype} \\
        {\bf typestring()} \autorefph{sec:iftypestring} \\
        {\bf primtypestring()} \autorefph{sec:ifprimtypestring} \\
        {\bf imtotargtype()} \autorefph{sec:ifimtotargtype} \\
        {\bf typeval()} \autorefph{sec:iftypeval} \\
        {\bf typevalstring()} \autorefph{sec:iftypevalstr} \\
        {\bf typeisequal()} \autorefph{sec:iftypeiseq} \\
 

    \section{time - get the current millisecond time}\label{sec:iftime}
        
        {\bf Call:} \\
        time(start)
        
        {\bf parameters:} \\
        No arguments
        
        {\bf Parameter=argument allowed:} \\
        No
        
        {\bf return mode:} \\
        immediate

        {\bf return type:} \\
        int
        
        {\bf Description:} \\
        Returns the current millisecond time, i.e. the time in millisecond
        since midnight, January 1, 1970 UTC.

        See also - \\
        \autorefph{sec:iftimestring} {\bf timestring()} for date and time
        as a string,
        \autorefph{sec:ifelapsedtime} {\bf elapsedtime()} for
        elapsed time.
 

    \section{timestring - get the date and time as a string}\label{sec:iftimestring}
        
        {\bf Call:} \\
        timestring(fmt)
        
        {\bf parameters:} \\
        \begin{tabular}{| l | l | l | l |}
        \hline
        param  & mode & type & \\
        \hline \hline
        fmt & I & str & optional format string   \\ \hline
        \end{tabular}
        
        {\bf Parameter=argument allowed:} \\
        No
        
        {\bf return mode:} \\
        immediate

        {\bf return type:} \\
        "str"
        
        {\bf Description:} \\
        Returns a string giving the current date and time. If the optional
        argument is supplied it must be a format string for the date/time.
        The format required is described below. The default is
	{\bf "yyyy-MM-dd HH:mm:ss Z"}.

        The following description of the optional format argument is taken
        from the Java API.
        
        Date and time formats are specified by date and time pattern strings.
        Within date and time pattern strings, unquoted letters from 'A' to 'Z'
        and from 'a' to 'z' are interpreted as pattern letters representing the
        components of a date or time string. Text can be quoted using single
        quotes (') to avoid interpretation. "''" represents a single quote. All
        other characters are not interpreted; they're simply copied into the
        output string during formatting or matched against the input string
        during parsing.

        The following pattern letters are defined (all other characters from 'A'
        to 'Z' and from 'a' to 'z' are reserved):

        \begin{tabular}{| l | l | l | l |}
        \hline
        param  & mode & type & \\
        \hline \hline
        fmt & I & str & optional format string   \\ \hline
        \end{tabular}


        \begin{tabular}{| l | l | l | l |}
        \hline
        Letter	&Date or Time Component	&Presentation	&Examples \\
        \hline \hline
        G	&Era designator	       &Text	       &AD   \\ \hline
        y	&Year	               &Year           &1996; 96   \\ \hline
        M	&Month in year	       &Month          &July; Jul; 07   \\ \hline
        w	&Week in year	       &Number         &27   \\ \hline
        W	&Week in month	       &Number         &2   \\ \hline
        D	&Day in year	       &Number         &189   \\ \hline
        d	&Day in month	       &Number         &10   \\ \hline
        F	&Day of week in month  &Number         &2   \\ \hline
        E	&Day in week	       &Text           &Tuesday; Tue   \\ \hline
        a	&Am/pm marker	       &Text           &PM   \\ \hline
        H	&Hour in day (0-23)    &Number         &0   \\ \hline
        k	&Hour in day (1-24)    &Number         &24   \\ \hline
        K	&Hour in am/pm (0-11)  &Number         &0   \\ \hline
        h	&Hour in am/pm (1-12)  &Number         &12   \\ \hline
        m	&Minute in hour	       &Number         &30   \\ \hline
        s	&Second in minute      &Number         &55   \\ \hline
        S	&Millisecond	       &Number         &978   \\ \hline
        \end{tabular}

        Pattern letters are usually repeated, as their number determines the
        exact presentation:

        Text: For formatting, if the number of pattern letters is 4 or more, the
        full form is used; otherwise a short or abbreviated form is used if
        available. For parsing, both forms are accepted, independent of the
        number of pattern letters.

        Number: For formatting, the number of pattern letters is the minimum
        number of digits, and shorter numbers are zero-padded to this amount.
        For parsing, the number of pattern letters is ignored unless it's needed
        to separate two adjacent fields.

        Year: If the formatter's Calendar is the Gregorian calendar, the
        following rules are applied. For formatting, if the number of pattern
        letters is 2, the year is truncated to 2 digits; otherwise it is
        interpreted as a number. For parsing, if the number of pattern letters
        is more than 2, the year is interpreted literally, regardless of the
        number of digits. So using the pattern "MM/dd/yyyy", "01/11/12" parses
        to Jan 11, 12 A.D.

        For parsing with the abbreviated year pattern ("y" or "yy"),
        SimpleDateFormat must interpret the abbreviated year relative to some
        century. It does this by adjusting dates to be within 80 years before
        and 20 years after the time the SimpleDateFormat instance is created.
        For example, using a pattern of "MM/dd/yy" and a SimpleDateFormat
        instance created on Jan 1, 1997, the string "01/11/12" would be
        interpreted as Jan 11, 2012 while the string "05/04/64" would be
        interpreted as May 4, 1964. During parsing, only strings consisting of
        exactly two digits, as defined by Character.isDigit(char), will be
        parsed into the default century. Any other numeric string, such as a one
        digit string, a three or more digit string, or a two digit string that
        isn't all digits (for example, "-1"), is interpreted literally. So
        "01/02/3" or "01/02/003" are parsed, using the same pattern, as Jan 2, 3
        AD. Likewise, "01/02/-3" is parsed as Jan 2, 4 BC. Otherwise, calendar
        system specific forms are applied. For both formatting and parsing, if
        the number of pattern letters is 4 or more, a calendar specific long
        form is used. Otherwise, a calendar specific short or abbreviated form
        is used. Month: If the number of pattern letters is 3 or more, the month
        is interpreted as text; otherwise, it is interpreted as a number.

        
        \autorefph{sec:iftime} {\bf time()} for date and time elements, \\
        \autorefph{sec:ifelapsedtime} {\bf elapsedtime()} for elapsed time.
        


    \section{toarray - convert to a one-dimensional array}\label{sec:iftoarray}
            
        {\bf Call:} \\
        toarray(v, n)
        
        {\bf parameters:} \\
        \begin{tabular}{| l | l | l | l |}
        \hline
        param  & mode & type & \\
        \hline \hline
        v & I &  -  & \\ \hline
        n & I & int & \\ \hline
        \end{tabular}
        
        {\bf Parameter=argument allowed:} \\
        No
        
        {\bf return mode:} \\
        immediate

        {\bf return type:} \\
        array
        
        {\bf Description:} \\
        Converts the 1st argument to a one-dimensional array.
        The argument must be a primitive type or a compound type where
        all members are of the same primitive type. Thus a primitive would be
        converted to an array of dimension 1. A compound type would be converted
        to a 1-dimensional array whose dimension equals the number of elements
        in the compound type.
        
        If the optional 2nd argument is provided it specifies the array size.
        Where the 1st argument is primitive or an array of dimension 1 it
        will be expanded to a one dimensional array of dimension n.
        For other types n must match the size of the compound type to be
        converted.
        


    \section{todegree - convert radian to degree}\label{sec:iftodegree}
            
        {\bf Call:} \\
        todegree(f)
        
        {\bf parameters:} \\
        \begin{tabular}{| l | l | l | l |}
        \hline
        param  & mode & type & \\
        \hline \hline
        f & I & float & \\ \hline
        \end{tabular}
        
        {\bf Parameter=argument allowed:} \\
        No
        
        {\bf return mode:} \\
        immediate

        {\bf return type:} \\
        float
        
        {\bf Description:} \\
        Converts the immediate argument from radian to degree.


    \section{tolower - convert string to lower case}\label{sec:iftolower}
        
        {\bf Call:} \\
        tolower(s)
        
        {\bf parameters:} \\
        \begin{tabular}{| l | l | l | l |}
        \hline
        param  & mode & type & \\
        \hline \hline
        s & I & str & subject string   \\ \hline
        \end{tabular}
        
        {\bf Parameter=argument allowed:} \\
        No
        
        {\bf return mode:} \\
        immediate

        {\bf return type:} \\
        str
        
        {\bf Description:} \\
        Returns the argument string with any upper case characters converted
        to lower case.


    \section{toradian - convert degree to radian}\label{sec:iftoradian}
            
        {\bf Call:} \\
        toradian()
        
        {\bf parameters:} \\
        \begin{tabular}{| l | l | l | l |}
        \hline
        param  & mode & type & \\
        \hline \hline
        f & I & float & \\ \hline
        \end{tabular}
        
        {\bf Parameter=argument allowed:} \\
        No
        
        {\bf return mode:} \\
        immediate

        {\bf return type:} \\
        float
        
        {\bf Description:} \\
        Converts the immediate argument from degree to radian.
  

    \section{tostring - string value}\label{sec:iftostring}
        
        {\bf Call:} \\
        tostring(a)
        
        {\bf parameters:} \\
        \begin{tabular}{| l | l | l | l |}
        \hline
        param  & mode & type & \\
        \hline \hline
        a & I & any &  \\ \hline
        \end{tabular}
        
        {\bf Parameter=argument allowed:} \\
        No
        
        {\bf return mode:} \\
        immediate

        {\bf return type:} \\
        str
        
        {\bf Description:} \\
        Returns a string representation of the immediate mode argument.
        Argument types accepted are str, int, uint, float, log, bits or an
        enumerated type. In the case of an enumerated type the string value
        returned is the value identifier.


    \section{toupper - convert string to upper case}\label{sec:iftoupper}
        
        {\bf Call:} \\
        toupper(s)
        
        {\bf parameters:} \\
        \begin{tabular}{| l | l | l | l |}
        \hline
        param  & mode & type & \\
        \hline \hline
        s & I & str & subject string   \\ \hline
        \end{tabular}
        
        {\bf Parameter=argument allowed:} \\
        No
        
        {\bf return mode:} \\
        immediate

        {\bf return type:} \\
        str
        
        {\bf Description:} \\
        Returns the argument string with any lower case characters converted
        to upper case.


    \section{trunc - truncate float or fixed to nearest integer towards zero}\label{sec:iftrunc}
        
        
        {\bf Call:} \\
        trunc(f)
        
        {\bf parameters:} \\
        \begin{tabular}{| l | l | l | l |}
        \hline
        param  & mode & type & \\
        \hline \hline
        f & I & float & \\ \hline
          & V & fixed, ufixed & \\ \hline
        \end{tabular}
        
        {\bf Parameter=argument allowed:} \\
        No
        
        {\bf return mode:} \\
        immediate or value

        {\bf return type:} \\
        int or uint
        
        {\bf Description:} \\
        Truncates an immediate floating point or target fixed point argument to the
        nearest integer towards zero. If the argument is {\bf immediate} mode
        the returned value is also {\bf immediate} mode. If the argument is mode
        {\bf selectvalue}, {\bf value}, {\bf static} or {\bf queue} the returned
        value is {\bf value} mode. Other modes are not allowed.

        For immediate mode input type "int" is returned. For target modes
        the returned value is type "int:" or "uint:" for input types
        "fixed:" and "ufixed:" respectively. If the argument is type "int:"
        or "uint:" then it is simply returned as the function value.


    \section{type - type of a variable or expression}\label{sec:iftype}
        
        {\bf Call:} \\
        type(v\_or\_e)
        
        {\bf parameters:} \\
        \begin{tabular}{| l | l | l | l |}
        \hline
        param  & mode & type & \\
        \hline \hline
        v\_or\_e & - &  -  & variable or expression \\ \hline
        \end{tabular}
        
        {\bf Parameter=argument allowed:} \\
        No
        
        {\bf return mode:} \\
        immediate

        {\bf return type:} \\
        type
        
        {\bf Description:} \\
        Return the type of a variable or expression.
        
        Note that if the argument is a variable of type {\bf type} then
        type {\bf type} will be returned, not the type contained in the type
        variable. For that use inbuilt function {\bf typeval()}.
        
        
        {\bf type()} \autorefph{sec:iftype} \\
        {\bf arraytype()} \autorefph{sec:ifarraytype} \\
        {\bf typestring()} \autorefph{sec:iftypestring} \\
        {\bf primtypestring()} \autorefph{sec:ifprimtypestring} \\
        {\bf imtotargtype()} \autorefph{sec:ifimtotargtype} \\
        {\bf targtoimtype()} \autorefph{sec:iftargtoimtype} \\
        {\bf typeval()} \autorefph{sec:iftypeval} \\
        {\bf typevalstring()} \autorefph{sec:iftypevalstr} \\
        {\bf typeisequal()} \autorefph{sec:iftypeiseq} \\


    \section{typeisequal - compare two types for equality}\label{sec:iftypeiseq}
        
        {\bf Call:} \\
        typeisequal(t1, t2, strict)
        
        {\bf parameters:} \\
        \begin{tabular}{| l | l | l | l |}
        \hline
        param  & mode & type & \\
        \hline \hline
        t1     & - &  -    & variable or expression \\ \hline
        t2     & - &  -    & variable or expression \\ \hline
        strict & I &  log  & variable or expression (optional) \\ \hline
        \end{tabular}
        
        {\bf Parameter=argument allowed:} \\
        Yes
        
        {\bf return mode:} \\
        immediate

        {\bf return type:} \\
        log
        
        {\bf Description:} \\
        Compare the types of the first two arguments. These may be
        immediate or target mode and may be variables or expressions.
        If an argument is of type {\bf type} the type value will be used
        rather than the type of the variable or expression. The argument types
        may be immediate or target but must be the same mode except in the case
        of type {\bf log}.
        
        Types {\bf int} and {\bf uint} are treated as equivalent.
        
        If either of the first two arguments are type {\bf type} then they
        will be regarded as equal regardless of their values. To compare
        type variables they must be passed via inbuilt function {\bf typeval()}
        which will unwrap the contained type.
        
        If the third argument is provided and is true the comparison will be
        strict, i.e. the types must be identical, except that {\bf int} and {\bf
        uint} are treated as equivalent (but the number of bits must be the
        same).
        
        If the third argument is omitted or is false and the first and second
        arguments are primitive types then they may vary in bit widths and still
        be regarded as equal. Thus the non-strict comparison returns true if an
        expression of the second type could be assigned to a variable of the
        first type.
        
        Where the first two arguments are compound types the third argument
        is ignored and the comparison will be strict (including bit widths).
        
        
        {\bf type()} \autorefph{sec:iftype} \\
        {\bf arraytype()} \autorefph{sec:ifarraytype} \\
        {\bf typestring()} \autorefph{sec:iftypestring} \\
        {\bf primtypestring()} \autorefph{sec:ifprimtypestring} \\
        {\bf imtotargtype()} \autorefph{sec:ifimtotargtype} \\
        {\bf targtoimtype()} \autorefph{sec:iftargtoimtype} \\
        {\bf typeval()} \autorefph{sec:iftypeval} \\
        {\bf typevalstring()} \autorefph{sec:iftypevalstr} \\
        {\bf typeisequal()} \autorefph{sec:iftypeiseq} \\


    \section{typestring - type of a variable as a string}\label{sec:iftypestring}
        
        {\bf Call:} \\
        typestring(ident)
        
        {\bf parameters:} \\
        \begin{tabular}{| l | l | l | l |}
        \hline
        param  & mode & type & \\
        \hline \hline
        ident & - &  -  & variable identifier \\ \hline
        \end{tabular}
        
        {\bf Parameter=argument allowed:} \\
        No
        
        {\bf return mode:} \\
        immediate

        {\bf return type:} \\
        str
        
        {\bf Description:} \\
        Return the type of a variable or expression as a string description.
        The type string is the same as that used for variable creation.
        
        If the argument is a variable which has a struct type or an enum type a
        string description of the struct or enum will be returned - see
        \autorefph{sec:ssecstructs} for a description of the string specification
        for structs and \autorefph{sec:ssecenums} for a description of the string
        specification for enums.
        
        If the argument is a variable of type "type" then this function will
        simply return "type", not a string representing the nested type. In that
        case use the inbuilt function {\bf typevalstring()}
        \autorefph{sec:iftypevalstr}.
        
        Note that a type string can also be formed from a type by simply assigning
        the type variable to a string variable or initialising the string variable
        with a type variable or the inbuilt function {\bf type()}.
        
        
        {\bf type()} \autorefph{sec:iftype} \\
        {\bf arraytype()} \autorefph{sec:ifarraytype} \\
        {\bf typestring()} \autorefph{sec:iftypestring} \\
        {\bf primtypestring()} \autorefph{sec:ifprimtypestring} \\
        {\bf imtotargtype()} \autorefph{sec:ifimtotargtype} \\
        {\bf targtoimtype()} \autorefph{sec:iftargtoimtype} \\
        {\bf typeval()} \autorefph{sec:iftypeval} \\
        {\bf typevalstring()} \autorefph{sec:iftypevalstr} \\
        {\bf typeisequal()} \autorefph{sec:iftypeiseq} \\


    \section{typeval - value of a variable of type "type"}\label{sec:iftypeval}
        
        {\bf Call:} \\
        typeval(v)
        
        {\bf parameters:} \\
        \begin{tabular}{| l | l | l | l |}
        \hline
        param  & mode & type & \\
        \hline \hline
        v & - &  -  & variable \\ \hline
        \end{tabular}
        
        {\bf Parameter=argument allowed:} \\
        No
        
        {\bf return mode:} \\
        immediate

        {\bf return type:} \\
        type
        
        {\bf Description:} \\
        Return the value of a variable or expression which is of type "type".
        This is used where the type value of a variable or expression of type
        "type" is required.
        
        
        {\bf type()} \autorefph{sec:iftype} \\
        {\bf arraytype()} \autorefph{sec:ifarraytype} \\
        {\bf typestring()} \autorefph{sec:iftypestring} \\
        {\bf primtypestring()} \autorefph{sec:ifprimtypestring} \\
        {\bf imtotargtype()} \autorefph{sec:ifimtotargtype} \\
        {\bf targtoimtype()} \autorefph{sec:iftargtoimtype} \\
        {\bf typeval()} \autorefph{sec:iftypeval} \\
        {\bf typevalstring()} \autorefph{sec:iftypevalstr} \\
        {\bf typeisequal()} \autorefph{sec:iftypeiseq} \\


    \section{typevalstring - value of a variable of type "type" as a string}\label{sec:iftypevalstr}
        
        {\bf Call:} \\
        typevalstring(v)
        
        {\bf parameters:} \\
        \begin{tabular}{| l | l | l | l |}
        \hline
        param  & mode & type & \\
        \hline \hline
        v & - &  -  & variable \\ \hline
        \end{tabular}
        
        {\bf Parameter=argument allowed:} \\
        No
        
        {\bf return mode:} \\
        immediate

        {\bf return type:} \\
        string
        
        {\bf Description:} \\
        Return the value of a variable or expression which is of type "type"
        as a string.
        This is used where the type value of a variable or expression of type "type"
        is required as a string. For example the code -
        \begin{lstlisting}[gobble=8, language=threepl]
           struct(s,
               f1, "log",
               f2, "int"
           );
           println(typestring(s));
        \end{lstlisting}
        will print "type" since variable {\bf s} has type "type". To get the type
        contained in the variable use typeval(s), thus the code
        \begin{lstlisting}[gobble=8, language=threepl]
           struct(s,
               f1, "log",
               f2, "int"
           );
           println(typevalstring(s));
        \end{lstlisting}
        will print "(f1,log,f2,int)".
        
        The argument may be -
        \blist
        \item   a variable of type "type" - the type it represents will
                be returned as a string
        \item   a string - the string will be returned unchanged
        \item   a variable - the type of the variable will be returned
                as a string
        \blend
        
        
        {\bf type()} \autorefph{sec:iftype} \\
        {\bf arraytype()} \autorefph{sec:ifarraytype} \\
        {\bf typestring()} \autorefph{sec:iftypestring} \\
        {\bf primtypestring()} \autorefph{sec:ifprimtypestring} \\
        {\bf imtotargtype()} \autorefph{sec:ifimtotargtype} \\
        {\bf targtoimtype()} \autorefph{sec:iftargtoimtype} \\
        {\bf typeval()} \autorefph{sec:iftypeval} \\
        {\bf typevalstring()} \autorefph{sec:iftypevalstr} \\
        {\bf typeisequal()} \autorefph{sec:iftypeiseq} \\


    \section{used - target value which is "true" when a function executes}\label{sec:ifused}
        
        {\bf Call:} \\
        used()
        
        {\bf parameters:} \\
        None.
        
        {\bf return mode:} \\
        value

        {\bf return type:} \\
        log
        
        {\bf Description:} \\
        This function returns a target signal of type "log" which is true
        when the function code actually executes (is evaluated).
        
        The following example uses a function which returns successive values
        in a sequence. It should advance the sequence only when the target
        assignment in which it is called executes.
        \begin{lstlisting}[gobble=8, language=threepl]
           seq while (..) {
               .
               a <- s();
               .
           }
           .
           .
           .
           function
           s () {
               var(v, "[5]int", {5, 7, 9, 3, 1});
               value(ret, "int:4");
               value(step, "log", used());
               del(ret <- ret, dimension(v), v, step);
               return(ret);
           }
        \end{lstlisting}
        Here the call to {\bf del()} creates a module which provides the
        sequence (since it is a module it does not constitute executable
        statements within the function definition - a function cannot contain
        inline target statements). The module steps to the next member of the
        sequence when its last argument {\bf step} is true. Here that argument
        reflects the inbuilt function {\bf used()} which will only be true on
        clock cycles on which the function value is being used (an intermediate
        value variable has been used here only to make the code more readable
        since passing {\bf used()} as an argument to {\bf del()} would
        not indicate the purpose of that argument).
        
        Note that the same instanciation of a function call can be used
        in more than one statement, for example -
        \begin{lstlisting}[gobble=8, language=threepl]
           value(sv, "int:4", s());
           seq while (..) {
               .
               a <- sv;
               .
               b <- sv;
           }
        \end{lstlisting}
        where the two assignments will get successive members of the sequence,
        which is quite different to -
        \begin{lstlisting}[gobble=8, language=threepl]
           seq while (..) {
               .
               a <- s();
               .
               b <- s();
           }
        \end{lstlisting}
        where each function call is a separate (duplicate) instance and hence
        they operate independently and produce the same values.
        
        see also inbuilt procedurem{\bf assert()} \autorefph{sec:ipassert}.


    \section{varexists - determine if a variable exists}\label{sec:ifvarexists}
        
        {\bf Call:} \\
        varexists(ident)
        
        {\bf parameters:} \\
        \begin{tabular}{| l | l | l | l |}
        \hline
        param  & mode & type & \\
        \hline \hline
        ident & - &  -  & identifier \\ \hline
        \end{tabular}
        
        {\bf Parameter=argument allowed:} \\
        No
        
        {\bf return mode:} \\
        immediate

        {\bf return type:} \\
        log
        
        {\bf Description:} \\
        Search currently accessible scopes and determine if a variable matching the
        identifier argument exists. If the identifier has {\bf global.}, {\bf
        file.}, {\bf calling.} or {\bf local.} prepended then the search will be
        restricted to the specified scope.
        
        To test for the existence of a directive
        see {\bf direxists()} \autorefph{sec:ifdirexists}.
        

        


    \section{varlist - return a list of target mode variables}\label{sec:ifvarlist}
        
        {\bf Call:} \\
        varlist(o)
        
        {\bf parameters:} \\
        \begin{tabular}{| l | l | l | l |}
        \hline
        param  & mode & type & \\
        \hline \hline
        o & I &  log  & selection option \\ \hline
        \end{tabular}
        
        {\bf Parameter=argument allowed:} \\
        No
        
        {\bf return mode:} \\
        immediate

        {\bf return type:} \\
        list
        
        {\bf Description:} \\
        This function returns a list of the currently declared target mode variables.
        If option {\bf o} is true then only INPUT, OUTPUT and CLOCK mode variables
        are included in the list.
        
        This function is of use when generating constraint files at the end of
        immediate execution. Attributes can be read from INPUT, OUTPUT and CLOCK mode
        variables and used to generate the constraints. Library procedure {\bf
        generateConstraints()} \autorefph{sec:lpgenerateconstraints} in
        xilinx/common.3pl does this in the case of Xilinx platforms, the call being
        placed last in trailer code.


    \section{wasassigned - determine if variable has been assigned}\label{sec:ifwasassigned}
        
        
        {\bf Call:} \\
        wasassigned(v)
        
        {\bf parameters:} \\
        \begin{tabular}{| l | l | l | l |}
        \hline
        param  & mode & type & \\
        \hline \hline
        v & - & - & \\ \hline
        \end{tabular}
        
        {\bf Parameter=argument allowed:} \\
        No
        
        {\bf return mode:} \\
        immediate

        {\bf return type:} \\
        log
        
        {\bf Description:} \\
        Returns true if the argument variable has been assigned. Assignment does
        not include initialisation. The argument must be a variable, can be of
        any type and can be any mode except INPUT, CMEMORY or RMEMORY.
        

    \section{waswritten - determine when a static or queue variable has been written}\label{sec:ifwaswritten}
        
        {\bf Call:} \\
        waswritten(v)
        
        {\bf parameters:} \\
        \begin{tabular}{| l | l | l | l |}
        \hline
        param  & mode & type & \\
        \hline \hline
        v & S Q & any & static or queue mode variable \\ \hline
        \end{tabular}
        
        {\bf Parameter=argument allowed:} \\
        No
        
        {\bf return mode:} \\
        value

        {\bf return type:} \\
        log
        
        {\bf Description:} \\
        Returns a target value of type "log" which is true on the clock
        cycle following a write to a static or queue variable or a part
        of a static variable. The argument must be a static or queue variable.
        
        A static variable argument may have struct fields or array subscripts in
        which case the returned value will be true following a write to any
        field or array member within those specified by the argument.
        
        A queue variable argument must not have subscripts or fields as a write
        to a queue variable always writes all components of that variable so
        it would be meaningless to specify some components of the queue variable.
        
        The argument variable write domain must be the same as the clock domain
        of the call to {\bf iswritten()}. If a static or queue mode variable does
        not yet have a (write) clock domain established then the call to
        {\bf iswritten()} will establish that domain.

        See also inbuilt function \\
        {\bf iswritten()} \autorefph{sec:ifiswritten}
        to detect the cycle on which the write occurs. \\
        {\bf isread()} \autorefph{sec:ifisread} can be used to signal on the
        cycle a queue is read.




    \section{width - get the width of a primitive type - DEPRECATED!}\label{sec:ifwidth}
        
        {\bf Call:} \\
        width(value)
        
        {\bf parameters:} \\
        \begin{tabular}{| l | l | l | l |}
        \hline
        param  & mode & type & \\
        \hline \hline
        value & I V SV S Q PR & -   & variable or expression   \\ \hline
        \end{tabular}
        
        {\bf Parameter=argument allowed:} \\
        No
        
        {\bf return mode:} \\
        immediate

        {\bf return type:} \\
        int
        
        {\bf Description:} \\
        Returns the width in bits of an immediate or target primitive type. The
        argument can be an expression or a variable. If the type is not a
        primitive or the mode is memory or clock a fatal error occurs. For
        an immediate argument the type must be int or log.
        For a +ve integer the sign bit is not included.
        
        This function has been deprecated. The inbuilt function {\bf size()}
        \autorefph{sec:ifsize} should be used instead.





    \section{words - get the size in words}\label{sec:ifwords}
        
        {\bf Call:} \\
        words(v)
        
        {\bf parameters:} \\
        \begin{tabular}{| l | l | l | l |}
        \hline
        param  & mode & type & \\
        \hline \hline
        v & V SV S Q PR & -   & variable \\ \hline
        \end{tabular}
        
        {\bf Parameter=argument allowed:} \\
        No
        
        {\bf return mode:} \\
        immediate

        {\bf return type:} \\
        int
        
        {\bf Description:} \\
        Returns the number of words.
        
        The argument cannot be cmemory, rmemory or clock mode.
        
        The function returns the number of words in the type of the variable, a
        word being a component of primitive type. The argument may have
        subscripts and fields in which case the number of words in the
        referenced part of the variable is returned.
        If the variable is type "null", 0 is returned.
        If the variable is type \verb'"->"', "list" or "map", 1 is returned
        (use inbuilt function {\bf entries()} \autorefph{sec:ifentries} for
        the number of entries in a list or map).
        If the variable is type "type", the number of words in the type
        is returned.
        


    \section{writedomain - get the write clock domain for a variable}\label{sec:ifwritedomain}
        
        {\bf Call:} \\
        writedomain(ident)
        
        {\bf parameters:} \\
        \begin{tabular}{| l | l | l | l |}
        \hline
        param  & mode & type & \\
        \hline \hline
        ident & - &  -  & variable identifier \\ \hline
        \end{tabular}
        
        {\bf Parameter=argument allowed:} \\
        No
        
        {\bf return mode:} \\
        clock

        {\bf return type:} \\
        none
        
        {\bf Description:} \\
        Returns the write (input or assignment) clock domain for a variable. If
        the variable has not been assigned so far null will be returned as the
        variable will not yet have had its input clock determined. The
        variable may be selectvalue, static, queue or priority mode.
        It cannot be memory or clock mode.
        
        See also - \\
        {\bf readdomain()} \autorefph{sec:ifreaddomain} which gets the read
        (output or evaluation) clock and \\
        {\bf domain()} \autorefph{sec:ifdomain} which
        gets the single clock for selectvalue, value, static and priority mode
        variables.
        
        For a queue the read and write clocks may be the same or may differ.
        For selectvalue, value and priority mode variables there is only
        one clock and all three inbuilt functions will return that clock.






